% =====================================================================
%  Quantum Error-Correcting Codes from Galois Hulls of Cyclic Codes
%  over Noncommutative Matrix Rings
%
%  Phase 8 / Task I.  Compile with:
%      pdflatex main.tex ; bibtex main ; pdflatex main.tex ; pdflatex main.tex
%
%  NOTE ON AUTHOR METADATA: no author names, affiliations, funding,
%  acknowledgements or declarations have been supplied.  The three
%  \AUTHORPLACEHOLDER marks below are the ONLY placeholders in this
%  manuscript and are deliberately visible.  See PAPER/README.md.
% =====================================================================
\documentclass[11pt,a4paper]{article}

\usepackage[margin=2.6cm]{geometry}
% NOTE: \usepackage[T1]{fontenc} is intentionally NOT used.  See
% compile/README.md: the recovered engine's texmf tree has no EC (T1)
% font metrics, so T1 silently falls back to nullfont.  OT1 (the
% article-class default) renders this manuscript identically: the .tex
% sources contain no non-ASCII characters.  On a current TeX Live,
% either setting compiles.
\usepackage[utf8]{inputenc}
\usepackage{amsmath,amssymb,amsthm}
\usepackage{array,booktabs,longtable}
\usepackage{enumitem}
\usepackage{xcolor}
\usepackage{url}
% File and directory names are long and contain underscores, so they are set
% with url.sty's \path, which may break at '_', '/' and '.'.  A plain
% \texttt{...} cannot break there and pushes into the margin.
\newcommand{\filepath}{\path}
% Typesetting: allow extra interword stretch on lines that would otherwise be
% pushed into the right margin.  This is a last-resort stretch (TeX's third
% pass) and is only used on lines that would otherwise be overfull; it never
% loosens a line that sets correctly.
\emergencystretch=2em
\usepackage[hypertexnames=false]{hyperref}
\hypersetup{colorlinks=true,linkcolor=blue!55!black,citecolor=green!35!black,
            urlcolor=blue!45!black,pdfauthor={\AUTHORPLACEHOLDER}}

% ------------------------------ notation ------------------------------
\newcommand{\F}{\mathbb F}
\newcommand{\Fq}{\mathbb F_q}
\newcommand{\At}{A_t}
\newcommand{\Rring}{R}
\newcommand{\Rn}{R_n}
\newcommand{\Bq}{B}
\newcommand{\bb}{\mathbf b}
\newcommand{\tr}{\operatorname{tr}}
\newcommand{\Tr}{\operatorname{Tr}}
\newcommand{\ann}{\operatorname{ann}}
\newcommand{\rev}{\operatorname{rev}}
\newcommand{\Hull}{\operatorname{Hull}}
\newcommand{\swt}{\operatorname{swt}}
\newcommand{\wt}{\operatorname{wt}}
\newcommand{\Aut}{\operatorname{Aut}}
\newcommand{\GL}{\operatorname{GL}}
\newcommand{\Jac}{\operatorname{Jac}}
\newcommand{\rad}{\operatorname{rad}}
\newcommand{\im}{\operatorname{im}}
\newcommand{\id}{\mathrm{id}}
\newcommand{\Nlen}{N}
\newcommand{\MU}{\mu}
\newcommand{\gperp}[1]{{}^{\perp_{#1}}}
\newcommand{\hperp}{\,{}^{\perp_{\widehat{}}}}
\newcommand{\hhatpers}{\,{}^{\perp^{\widehat{}}_s}}
\newcommand{\omegaperp}{{}^{\perp_\omega}}
\newcommand{\AUTHORPLACEHOLDER}{\textbf{[\,AUTHOR METADATA TO BE
  SUPPLIED\,]}}
\newcommand{\blanksymbol}{\textbf{[TO BE SUPPLIED]}}

% ------------------------------ theorems ------------------------------
\theoremstyle{plain}
\newtheorem{theorem}{Theorem}[section]
\newtheorem{proposition}[theorem]{Proposition}
\newtheorem{lemma}[theorem]{Lemma}
\newtheorem{corollary}[theorem]{Corollary}
\theoremstyle{definition}
\newtheorem{definition}[theorem]{Definition}
\newtheorem{example}[theorem]{Example}
\newtheorem{assumption}[theorem]{Assumption}
\theoremstyle{remark}
\newtheorem{remark}[theorem]{Remark}
\newtheorem{openproblem}[theorem]{Open Problem}
\newtheorem{conjecture}[theorem]{Conjecture}

% ------------------------------ metadata ------------------------------
\title{Quantum Error-Correcting Codes from Galois Hulls of Cyclic
       Codes over Noncommutative Matrix Rings}
\author{\AUTHORPLACEHOLDER}
\date{}

\begin{document}
\maketitle

\begin{abstract}
We construct $q$-ary quantum error-correcting codes from the Galois hulls of
cyclic codes over the noncommutative matrix ring $R=M_2(A_t)$,
$A_t=\F_q[u]/\langle u^t\rangle$.  The duality is driven by the
anti-automorphism $\theta_s=\sigma_s\circ(\cdot)^{\mathsf T}$, where
$\sigma_s$ is the coefficient Frobenius $a\mapsto a^{p^s}$; codes are right
ideals of $R_n=R[X]/\langle X^n-1\rangle$.  We prove the duality and hull
theory for arbitrary right ideals, with no principality hypothesis, and we
refute the natural conjecture that Galois hulls are principal: hulls with
minimal generator number $\MU_B=3$ exist, and we show that $\MU_B\ge3$ forces
$t\ge2$ and $p\mid n$.  The coefficient-expansion map to $\F_q^{4tn}$ is shown
to transport the Galois pairing to an explicitly identified
$\widehat\sigma_s$-Hermitian form with Hankel Gram matrix $I_{4n}\otimes H_t$
--- not the Euclidean form unless $t=1$ --- and to preserve both duality and
hulls, including the non-principal ones.  Every Galois hull is self-orthogonal
for the transported form, with no reflexivity hypothesis.  This yields three
quantum routes of length $\Nlen=4tn$: a Galois-native twisted stabilizer route,
a Euclidean/CSS route, and an entanglement-assisted route
$[[N,N-k-h_E,d;k-h_E]]_q$.  We prove that the distance is the minimum weight of
the quotient code $D^{\perp_E}/\Hull_E(D)$, and derive from it a coset
Singleton bound $K+d\le N+1$ that is neither the entanglement-assisted
Singleton bound nor implied by it.  Exhaustive computation over the enlarged
family list certifies $28$ codes by three independent methods, including the
first stabilizer distance $4$ in this framework, $[[24,8,4]]_9$ and
$[[20,4,4]]_4$, and entanglement-assisted codes up to $[[20,2,10;18]]_4$.  No
record parameters are claimed.

\smallskip
\noindent\textbf{Keywords.} cyclic codes over rings; matrix rings; Galois
duality; Galois hulls; Gray map; quantum stabilizer codes;
entanglement-assisted quantum codes.

\smallskip
\noindent\textbf{MSC 2020.} 94B05, 94B15, 94B60, 16S50, 11T71.
\end{abstract}

\tableofcontents
\bigskip

% ------------------------------ sections ------------------------------
% =====================================================================
\section{Introduction and research motivation}
\label{sec:intro}
% =====================================================================

The construction of quantum error-correcting codes from classical codes over
rings is by now a mature enterprise, and one of its most productive branches
runs through the \emph{hull} of a code: for a code $C$ equipped with a pairing
$\langle\cdot,\cdot\rangle$, the hull $C\cap C^{\perp}$ measures how far $C$ is
from being self-orthogonal, and it is precisely the hull dimension that
determines the parameters of the quantum code one can extract. Euclidean hulls
give CSS codes, Hermitian hulls give $q$-ary stabilizer codes, and --- when the
hull is small rather than large --- the entanglement-assisted formalism converts
\emph{any} classical code into a quantum code at the price of shared
entanglement.

Parallel to this, the Galois (or $\ell$-Galois) hull has been introduced over
finite fields \cite{FanZhang2017GaloisSelfDual,LiuPan2020GaloisHulls} and used
to produce entanglement-assisted codes with good parameters
\cite{Cao2021MDSGaloisHulls}. The Galois parameter $\ell$ interpolates between
the Euclidean and the Hermitian inner product, and because it can be tuned it
gives a strictly larger supply of codes with a prescribed hull dimension.

What has been missing is the passage to a \emph{non-commutative} alphabet. There
is a well-developed theory of cyclic codes over matrix rings
\cite{Alahmadi2013CyclicM2F2,LuoParampalli2018M2F2uF2,Dinh2024CyclicRingsOfMatrices,
JoshiPatelPrakash2026M4}, and there is a well-developed theory of quantum codes
from Galois hulls over commutative rings, culminating in the recent
constacyclic construction over a finite non-chain ring
\cite{ZhangKongZheng2026QuantumGaloisHulls}. The two have not met, and they do
not meet trivially: over a commutative ring the Galois twist is an
automorphism applied to one argument of the inner product, whereas over a matrix
ring the twist must be composed with the transpose to produce an
\emph{anti}-automorphism, and only then does the resulting pairing give a
perfect duality on \emph{right} modules. This paper builds that bridge and
constructs quantum codes at the far end of it.

\subsection{The chain of constructions}

Throughout, $q=p^r$, and
\[
  A_t=\F_q[u]/\langle u^t\rangle,\qquad
  R=M_2(A_t),\qquad
  R_n=R[X]/\langle X^n-1\rangle,\qquad
  \Nlen=4tn .
\]
A \emph{cyclic code of length $n$ over $R$} is a right ideal of $R_n$; this is
the right-module convention fixed for this alphabet in
\cite{Alahmadi2013CyclicM2F2}, and Lemma~\ref{lem:failure} below shows that it
is not a cosmetic choice. Under the Morita identification $R_n\cong M_2(B)$ with
$B=A_t[X]/\langle X^n-1\rangle$ \emph{commutative}, right ideals of $R_n$
correspond to $B$-submodules of $B^2$.

The contribution of this paper is the following chain, each link of which is
proved here:
\[
  \begin{gathered}
    \boxed{\text{noncommutative matrix ring}}
    \;\longrightarrow\;
    \boxed{\text{cyclic codes}}
    \;\longrightarrow\;
    \boxed{\text{Galois duality and hulls}}\\[3pt]
    \longrightarrow\;
    \boxed{\text{Gray map}}
    \;\longrightarrow\;
    \boxed{\text{quantum codes}} .
  \end{gathered}
\]
In more detail:

\begin{enumerate}[label=\textnormal{(\arabic*)},leftmargin=*]
\item \textbf{Galois duality on a non-commutative ring.} With
  $\sigma_s(a)=a^{p^s}$ the coefficient Frobenius (fixing $u$ and $X$) and
  $\theta_s=\sigma_s\circ(\cdot)^{\mathsf T}$, the pairing
  $\bb_s(x,y)=\Lambda\bigl(\sum_i\theta_s(x_i)y_i\bigr)$, where
  $\Lambda=\lambda\circ\tr$ and $\lambda$ extracts the coefficient of
  $u^{t-1}$, is non-degenerate on right $R_n$-modules, and
  $\dim C+\dim C^{\perp_s}=\Nlen$ for every code (Theorem~\ref{thm:duality}).
  We give the dual explicitly for \emph{arbitrary} finitely generated right
  ideals (Theorem~\ref{thm:dualformula}), including the reversal operator,
  which is load-bearing exactly when $n\ge3$.

\item \textbf{Hulls, with no principality hypothesis.} We prove the structural
  properties of $\Hull_s(C)=C\cap C^{\perp_s}$ for arbitrary right ideals, and
  we \emph{refute} the natural conjecture that a Galois hull is principal: hulls
  with minimal generator number $\MU_B=3$ exist
  (Theorem~\ref{thm:nonprincipal}). The phenomenon is confined:
  $\MU_B\ge3$ forces $t\ge2$ and $p\mid n$.

\item \textbf{The exact induced form.} The coefficient expansion
  $\varphi_n:R^n\to\F_q^{\Nlen}$ is not an isometry in any interesting sense,
  but it \emph{is} a form-isometry for exactly one form, and we identify that
  form: $\langle u,v\rangle^{\widehat{}}_s=u^{[s]}Jv^{\mathsf T}$ with
  $J=I_{4n}\otimes H_t$ and $H_t$ the Hankel reversal matrix
  (Theorem~\ref{thm:gram}). Duality and hulls are transported exactly, and the
  transport is valid on the non-principal hulls as well
  (Theorem~\ref{thm:grayhull}).

\item \textbf{Every Galois hull is self-orthogonal} for the transported form,
  by a one-sided argument that uses neither reflexivity nor principality
  (Theorem~\ref{thm:hullso}). This is the hinge of the whole construction: it
  is what makes the quantum step available in the general Galois regime rather
  than only in the reflexive one.

\item \textbf{Quantum codes of length $\Nlen=4tn$.} Three routes are available
  and all three are proved from scratch in symplectic language: a Galois-native
  twisted stabilizer route R1, a Euclidean/CSS route R2 and an
  entanglement-assisted route R3 giving
  $[[\Nlen,\Nlen-k-h_E,d;k-h_E]]_q$ (Theorem~\ref{thm:parameters}).

\item \textbf{A coset-distance theorem and a new bound.} The distance is the
  minimum weight of the \emph{quotient} code
  $D^{\perp_E}/\Hull_E(D)$, of dimension $K$ and length $\Nlen$
  (Theorem~\ref{thm:cosetdistance}); puncturing $\Nlen-d+1$ coordinates is
  injective on that quotient, which gives the \emph{coset Singleton bound}
  $K+d\le \Nlen+1$ (Theorem~\ref{thm:cosetsingleton}). The latter is neither
  the entanglement-assisted Singleton bound nor implied by it.
\end{enumerate}

\subsection{What the paper does and does not claim}

The codes constructed here are new in \emph{origin}, not in parameters.
Exhaustive computation over the accessible family list ($q\in\{4,8,9\}$,
$tn\le8$) certifies $28$ codes by three independent methods, among them the
first stabilizer distance $4$ produced by this framework, $[[24,8,4]]_9$ and
$[[20,4,4]]_4$, and entanglement-assisted codes up to $[[20,2,10;18]]_4$. Every
stabilizer code we construct whose parameters can be compared with
\cite{GrasslCodetables} is at best equal to the best known distance --- and then
only with $d=2$ --- and every code with $d\ge3$ is one or two below it. We claim
no record. What we do claim is that the chain above is proved in the
non-commutative setting, that it produces quantum codes from hulls that are
\emph{not} principal, and that the two new bounds of item (6) hold.

\subsection{Organisation}

Section~\ref{sec:lit} reviews the literature and states the research gap.
Section~\ref{sec:prelim} fixes the rings, the right-module convention and the
anti-automorphism, and proves the failure lemma that motivates the convention.
Section~\ref{sec:cyclic} develops cyclic codes and Galois duality.
Section~\ref{sec:hulls} is the algebraic heart: hull structure, the
non-principal counterexample and its confinement.
Section~\ref{sec:gray} transports everything to $\F_q^{\Nlen}$ and identifies
the induced form. Section~\ref{sec:quantum} builds the quantum codes.
Section~\ref{sec:cert} proves the coset-distance theorem and the coset Singleton
bound, describes the certification methodology and reports the parameter
analysis. Section~\ref{sec:cmp} compares with the literature and states the
novelty precisely. Section~\ref{sec:conclusion} collects open problems. Appendix~A
contains the full parameter tables and Appendix~B the reproducibility package.

% =====================================================================
\section{Literature review and the research gap}
\label{sec:lit}
% =====================================================================

\subsection{Cyclic codes over matrix rings}

The use of a matrix ring as a code alphabet goes back to
Bachoc's construction of modular lattices, and the coding-theoretic study of
cyclic codes over $M_2(\F_2)$ was initiated by Alahmadi, Sboui, Sol\'e and
Yemen \cite{Alahmadi2013CyclicM2F2}, who fixed the convention that cyclic codes
are \emph{right} ideals and used the ring-valued Euclidean dot product. Luo and
Parampalli \cite{LuoParampalli2018M2F2uF2} extended the structure theorem to
$M_2(\F_2+u\F_2)$, i.e. to $t=2$ in our notation, and Dinh, Gaur, Kumar, Singh
and Singh \cite{Dinh2024CyclicRingsOfMatrices} to $M_2(\F_2[u]/\langle u^k\rangle)$
for all $k$ --- the same ring family as the present paper --- obtaining the
structure via an isometry onto $\F_4[u,v]/\langle u^k,v^2\rangle$. Most
recently Joshi, Patel and Prakash \cite{JoshiPatelPrakash2026M4} treated
$M_4(\F_2[u]/\langle u^k\rangle)$ with Euclidean and Hermitian duals and both
the Bachoc and the Gray map into $\F_{16}$. A separate, genuinely different line
uses skew polynomial rings and cyclic algebras over $M_2(\F_2)$
\cite{SiNiu2023SkewM2F2}.

What \emph{none} of these papers contains is a Galois duality, a hull, or a
quantum construction. Their duals are the Euclidean and the Hermitian dual, and
their output is a classical code over a finite field, compared with the best
known classical codes.

\subsection{Duality over finite rings}

The authority for duality over non-commutative finite rings is Wood
\cite{Wood1999Duality}: the extension of the MacWilliams equivalence theorem to
finite rings characterises Frobenius rings, and duality for modules over a
finite Frobenius ring is the correct general framework. We import the statement
that $R=M_m(A_t)$ is Frobenius and prove directly, in
Lemma~\ref{lem:lambda}, that $\Lambda=\lambda\circ\tr$ is a Frobenius functional
on it.

For the weight side, Greferath and Schmidt \cite{GreferathSchmidt1999GrayIsometry}
construct Gray isometries for finite \emph{chain} rings. Their theorem does not
apply to $M_2(A_t)$, which is not a chain ring for $t\ge1$ and $m=2$; we use it
only to delimit what our own map can achieve (Remark~\ref{rem:noisometry}).

\subsection{Galois duality and Galois hulls}

Over a finite field, Fan and Zhang \cite{FanZhang2017GaloisSelfDual} introduced
the $\ell$-Galois inner product $\langle a,b\rangle_\ell=\sum_i a_ib_i^{p^\ell}$
and studied $\ell$-Galois self-dual constacyclic codes; Liu and Pan
\cite{LiuPan2020GaloisHulls} introduced the $\ell$-Galois hull
$\Hull_\ell(C)=C\cap C^{\perp_\ell}$ and gave its dimension formulas. Cao
\cite{Cao2021MDSGaloisHulls} constructed MDS codes with Galois hulls of
arbitrary dimensions and the corresponding entanglement-assisted quantum codes.
All of this takes place over a \emph{commutative} base.

The closest prior work to the present paper is Zhang, Kong and Zheng
\cite{ZhangKongZheng2026QuantumGaloisHulls}, who study constacyclic codes over a
finite non-chain ring $R$ (commutative, with a decomposition
$C=e_1C_1\oplus\cdots\oplus e_tC_t$), determine the dimensions of the
$\ell$-Galois hulls of their Gray images, and construct EAQECCs and QECCs by a
construction $X$ for the Galois inner product, obtaining some QMDS codes and
some QECCs better than the best known ones.

\subsection{The research gap}

The gap is the conjunction of three properties, none of which is present in the
literature:

\begin{enumerate}[label=\textnormal{(\arabic*)},leftmargin=*]
\item the alphabet is \textbf{non-commutative} (a matrix ring over a finite
  chain ring), so that the Galois twist must be an \emph{anti}-automorphism and
  the left/right distinction is substantive;
\item the duality is \textbf{Galois} (a Frobenius parameter $s$ is present,
  interpolating between the Euclidean and Hermitian cases in a way that a
  commutative alphabet cannot reproduce);
\item the outcome is \textbf{quantum}: stabilizer and entanglement-assisted
  codes, obtained from the hull through a Gray map whose induced form is
  identified exactly.
\end{enumerate}

Property (1) alone is the matrix-ring literature; (2) alone is the Galois-hull
literature; (2)+(3) over a commutative non-chain ring is
\cite{ZhangKongZheng2026QuantumGaloisHulls}. The conjunction (1)+(2)+(3) is
what this paper supplies.

Two consequences of the non-commutativity are worth signalling here, because
they have no commutative analogue and they shape the whole paper.

\begin{itemize}
\item The pairing must use $\theta_s=\sigma_s\circ(\cdot)^{\mathsf T}$, not
  $\sigma_s$. Over a commutative ring the transpose is invisible and the
  $\ell$-Galois form is simply Frobenius-in-one-argument; here the transpose is
  what makes $\theta_s$ an anti-automorphism, and the anti-automorphism is what
  makes the dual of a \emph{right} module again a right module
  (Theorem~\ref{thm:duality}). The plain product $\sum_ix_iy_i$, which is the
  Euclidean dot product used in \cite{Alahmadi2013CyclicM2F2}, has a perfect
  dual on left modules but \emph{not} on right modules
  (Lemma~\ref{lem:failure}).
\item The hull need not be principal, so no generator-polynomial formula for
  the hull can exist, and every downstream statement must be proved without a
  principality hypothesis. This is not a technical nuisance: the
  counterexamples are exactly the families that the quantum constructions later
  use (Theorem~\ref{thm:nonprincipal} and Section~\ref{subsec:witnesses}).
\end{itemize}

\subsection{Quantum codes from classical codes}

The stabilizer formalism for $q$-ary quantum codes is due to
\cite{KetkarEtAl2006NonbinaryStabilizer} in the non-binary case. We do not
import the parameter formulas: Section~\ref{sec:quantum} re-derives the
stabilizer lemma and all three routes from the symplectic commutation relation,
so that every hypothesis (in particular the requirement that commutation be
tested on an $\F_p$-basis, Theorem~\ref{thm:fpbasis}) is visible where it is
used.

% =====================================================================
\section{Preliminaries and the noncommutative matrix-ring framework}
\label{sec:prelim}
% =====================================================================

\subsection{Standing hypotheses and notation}

Throughout, $p$ is a prime, $q=p^r$, and $t,n\ge1$ are integers. We set
\begin{equation}\label{eq:rings}
  A_t=\F_q[u]/\langle u^t\rangle,\qquad
  R=M_2(A_t),\qquad
  \Bq=A_t[X]/\langle X^n-1\rangle,\qquad
  \Rn=R[X]/\langle X^n-1\rangle ,
\end{equation}
and write $\Nlen=4tn$. Elements of $A_t$ are written $\sum_{a=0}^{t-1}c_au^a$.
All dimensions and all cardinalities of codes are over $\F_q$ unless a different
field is named explicitly; this is maintained without exception, and the only
place where $\F_p$ enters is Theorem~\ref{thm:fpbasis}.

\begin{definition}[cyclic code]\label{def:code}
A \emph{cyclic code of length $n$ over $R$} is a right ideal of $\Rn$. We
identify $\Rn$ with $R^n$ by $x=\sum_ex_eX^e\leftrightarrow(x_0,\dots,x_{n-1})$
and write $C\le \Rn$ for a right $\Rn$-submodule.
\end{definition}

The choice of right ideals is forced, not conventional: see
Lemma~\ref{lem:failure}. The ring $\Bq$ is commutative (it is a quotient of the
commutative ring $A_t[X]$), and there is a canonical isomorphism
\begin{equation}\label{eq:morita}
  \Rn \;\cong\; M_2(\Bq)
\end{equation}
because $X$ is central, so that a polynomial with matrix coefficients may be
read as a matrix of polynomials. Under \eqref{eq:morita}, right ideals of
$\Rn$ correspond to $\Bq$-submodules of $\Bq^2$; this is the Morita
correspondence we use repeatedly and we do not reprove it.

\begin{lemma}[basic structure]\label{lem:basic}
\begin{enumerate}[label=\textnormal{(\alph*)},leftmargin=*]
\item $A_t$ is a finite chain ring: its ideals are exactly
  $u^iA_t$, $0\le i\le t$; $|A_t|=q^t$; its maximal ideal is $uA_t$ and its
  residue field is $\F_q$.
\item $R$ is a finite Frobenius ring with $|R|=q^{4t}$,
  $\dim_{\F_q}R=4t$, $Z(R)=A_t\cdot I$, $\Jac(R)=M_2(uA_t)$, and exactly $t+1$
  two-sided ideals, namely $u^iR$.
\item $\dim_{\F_q}R^n=\Nlen=4tn$.
\end{enumerate}
\end{lemma}

\begin{proof}
(a) Every element of $A_t$ has a unique form $c_{i}u^i+\text{(higher terms)}$
with $c_i\ne0$, and then the element is $u^i$ times a unit; hence every ideal is
of the stated form and they form a chain. (b) The stated descriptions of the
centre and the Jacobson radical are immediate from the matrix-unit computation;
the two-sided ideals of $M_2(S)$ are $M_2(I)$ for ideals $I\trianglelefteq S$,
which gives $t+1$ of them. The Frobenius property is imported
\cite{Wood1999Duality}; a direct Frobenius functional is exhibited in
Lemma~\ref{lem:lambda} below. (c) is immediate from (b), since
$\dim_{\F_q}A_t=t$ and $\dim_{\F_q}R=4\dim_{\F_q}A_t$.
\end{proof}

\subsection{The Frobenius functional}\label{subsec:lambda}

Let $\lambda:A_t\to\F_q$ be the coefficient of $u^{t-1}$, and let
$\tr$ denote the matrix trace. Put
\begin{equation}\label{eq:Lambda}
  \Lambda=\lambda\circ\tr:R\longrightarrow\F_q,\qquad
  \Lambda(X)=\lambda(X_{11}+X_{22}) .
\end{equation}

\begin{lemma}\label{lem:lambda}
$\Lambda$ is symmetric ($\Lambda(XY)=\Lambda(YX)$), associative (the kernel
contains no nonzero left or right ideal), $\F_q$-linear, and non-degenerate:
the bilinear form $(X,Y)\mapsto\Lambda(XY)$ on $R$ is non-degenerate.
\end{lemma}

\begin{proof}
Symmetry: $\tr(XY)=\tr(YX)$ for matrices over any commutative ring, and $\lambda$
is $\F_q$-linear, so $\Lambda(XY)=\Lambda(YX)$.

Non-degeneracy: let $X=(X_{ij})\ne0$ and choose $(k,l)$ with $X_{kl}\ne0$;
write $X_{kl}=\sum_ac_au^a$ and let $a_0$ be the least index with $c_{a_0}\ne0$.
Take $Y$ to be the matrix whose $(l,k)$ entry is $u^{t-1-a_0}$ and whose other
entries are $0$. Then $XY$ has trace
$\sum_{i,j}X_{ij}Y_{ji}=X_{kl}u^{t-1-a_0}$, whose $u^{t-1}$-coefficient is
$c_{a_0}\ne0$. Hence $\Lambda(XY)\ne0$, and the form is non-degenerate.

Associativity (no nonzero one-sided ideal lies in $\ker\Lambda$): if
$X\ne0$, the $Y$ constructed above already gives $\Lambda(XY)\ne0$ with $Y$
arbitrary in $R$, so the right annihilator of $\Lambda$ in $R$ is zero; by
symmetry the same holds on the left.
\end{proof}

Together with the import that $R$ is Frobenius, Lemma~\ref{lem:lambda} says that
$(R,\Lambda)$ is a Frobenius algebra over $\F_q$, which is the setting we need.

\subsection{The right-module convention and the failure lemma}

\begin{lemma}[failure lemma]\label{lem:failure}
Let $S$ be a finite Frobenius ring and let
$C$ be a \textbf{left} ideal. Then the dual of $C$ with respect to the plain
product $\sum_ix_iy_i$ satisfies $|C|\cdot|C^{\perp}|=|S|^n$. For a
\textbf{right} ideal this fails in general. Concretely, for
$S=R=M_2(\F_2)$, $C=uR$ (with $t=2$, $q=2$, $n=1$, so that $u$ denotes a
nilpotent generator of $A_2$ embedded diagonally) one has
$|C|\cdot|C^{\perp}|=4\ne16=|R|$.
\end{lemma}

\begin{proof}
For a left ideal $C$, the map
$\psi_C:C\to\Hom_{\F_q}(S/C^{\perp},\F_q)$, $c\mapsto(y\mapsto\Lambda(cy))$, is
an $\F_q$-linear isomorphism by the non-degeneracy of $\Lambda$, and the
standard dimension count gives the size formula.

For right ideals the argument breaks: $\psi_C$ is no longer $\F_q$-balanced in
the way required, because the plain product is not compatible with the
right action. The stated counterexample is a direct computation: with
$R=M_2(\F_2)$ and $C=uR$ one has $|C|=4$ while the plain-product dual
$C^{\perp}=\{y:\tr(xy)=0\ \forall x\in C\}$ has size $1$, giving
$|C||C^{\perp}|=4\ne 16$.
\end{proof}

\begin{remark}\label{rem:convention}
Lemma~\ref{lem:failure} is the reason the pairing of
Definition~\ref{def:pairing} below carries the transpose: only then does the
dual of a right module come out as a right module with the correct size
(Theorem~\ref{thm:duality}). This is not a matter of notation.
\end{remark}

\subsection{Automorphisms and the anti-automorphism}

\begin{definition}\label{def:sigma}
For $s\ge0$ let $\sigma_s:\F_q\to\F_q$, $a\mapsto a^{p^s}$. Extend $\sigma_s$
to $A_t$, to $R$ and to $\Rn$ coefficient-wise, fixing $u$ and $X$. Define
\begin{equation}\label{eq:theta}
  \theta_s=\sigma_s\circ(\cdot)^{\mathsf T},
\end{equation}
the composite of the coefficient-wise $\sigma_s$ with the transpose.
\end{definition}

\begin{lemma}\label{lem:theta}
\begin{enumerate}[label=\textnormal{(\alph*)},leftmargin=*]
\item $\sigma_s\in\Aut(A_t)$ and its entrywise extension lies in $\Aut(R)$;
  $\sigma_s$ is nontrivial iff $r\nmid s$, and its order is $r/\gcd(r,s)$.
\item $\theta_s$ is an \textbf{anti}-automorphism of $R[X]$ and of $\Rn$:
  $\theta_s(ab)=\theta_s(b)\theta_s(a)$.
\item $\theta_s^2=\sigma_{2s}$; more generally $\theta_s^{2k}=\sigma_{2ks}$ and
  $\theta_s^{2k-1}=\theta_{(2k-1)s}$. In particular $\theta_s$ is an involution
  iff $r\mid 2s$.
\item $\theta_s(X^n-1)=X^n-1$, so $\theta_s$ is well defined on $\Rn$; it is
  additive, $\F_p$-linear and $\sigma_s$-semilinear over $\F_q$ (not
  $\F_q$-linear unless $r\mid s$).
\end{enumerate}
\end{lemma}

\begin{proof}
(a) $a\mapsto a^{p^s}$ is a field automorphism of $\F_q$ of order
$r/\gcd(r,s)$, and it fixes the coefficients of the relation $u^t=0$, so it
extends to $A_t$ by fixing $u$; entrywise extension preserves matrix
multiplication.

(b) For matrices over any ring, $(AB)^{\mathsf T}=B^{\mathsf T}A^{\mathsf T}$,
and $\sigma_s$ is a ring automorphism, hence
$\theta_s(ab)=\sigma_s((ab)^{\mathsf T})
=\sigma_s(b^{\mathsf T})\sigma_s(a^{\mathsf T})=\theta_s(b)\theta_s(a)$. The
same computation works in $R[X]$; note that there the index of $X$ is swapped
as well, which is why the statement is recorded separately for $R[X]$ and for
$\Rn$.

(c) $\theta_s^2(X)=\sigma_s(\sigma_s(X^{\mathsf T})^{\mathsf T})=\sigma_s^2(X)
=\sigma_{2s}(X)$; the general powers follow by induction.

(d) $\sigma_s$ fixes $X$ and the coefficients of $X^n-1$ are $0,\pm1\in\F_p$,
which $\sigma_s$ fixes. Semilinearity is immediate from $\sigma_s$ being a field
automorphism of order $>1$ in general.
\end{proof}

\begin{remark}
We use only the coefficient-Frobenius subgroup of $\Aut(A_t)$: an automorphism
of $A_t$ may move $u$, and we never use those. Everything in this paper
generalises verbatim to $\tau\in\Aut(A_t)$ with $\theta_\tau=\tau\circ(\cdot)^{\mathsf T}$
and $\theta_\tau^2=\tau^2$; we keep $\sigma_s$ for concreteness.
\end{remark}

\subsection{The CRT decomposition of $\Bq$}\label{subsec:crt}

Let $X^n-1=\prod_{j}g_j^{e_j}$ be the factorisation into pairwise coprime
monic irreducibles over $\F_q$. By the Chinese remainder theorem,
\begin{equation}\label{eq:crt}
  \Bq=A_t[X]/\langle X^n-1\rangle
  \;\cong\;\prod_j \Bq_j,\qquad
  \Bq_j=A_t[X]/\langle g_j^{e_j}\rangle ,
\end{equation}
and, writing $d_j=\deg g_j$, one has $\Bq_j\cong\F_{q^{d_j}}[u,w]/\langle u^t,w^{e_j}\rangle$.
The isomorphism \eqref{eq:crt} is realised explicitly as an $\F_q$-linear ring
isomorphism $\Psi$ (with matrix and inverse), so that it composes with the Gray
map, the Morita identification and the quantum routes.

\begin{lemma}[component generation]\label{lem:components}
Let $W\le \Bq^2$. Under \eqref{eq:crt}, $W=\prod_jW_j$ with
$W_j\le\Bq_j^2$, and
\[
  \MU_{\Bq}(W)=\max_j\MU_{\Bq_j}(W_j),
\]
where $\MU_S(W)$ is the minimal number of generators of $W$ as an
$S$-module. Moreover $\Bq_j$ is a finite chain ring iff $t=1$ or $e_j=1$; in
the remaining ``wild'' case ($t\ge2$ and $e_j\ge2$) the maximal ideal of
$\Bq_j$ is $\mathfrak m_j=(u,w)$, which is not principal, and the bound
$\MU_{\Bq_j}(W_j)\le2$ fails.
\end{lemma}

\begin{proof}
Generation is componentwise because a generating set may be chosen
componentwise and the idempotents of the product allow the components to be
addressed independently; this gives the displayed formula.

For the chain-ring criterion: $\Bq_j\cong\F_{q^{d_j}}[u,w]/\langle
u^t,w^{e_j}\rangle$ is local with maximal ideal $(u,w)$ when $t\ge2$ and
$e_j\ge2$; its maximal ideal then needs two generators, so the ring is not a
chain ring. If $t=1$ the ring is $\F_{q^{d_j}}[w]/\langle w^{e_j}\rangle$, a
chain ring; if $e_j=1$ it is $\F_{q^{d_j}}[u]/\langle u^t\rangle$, a chain
ring. Finally, a finite local ring is a chain ring iff its maximal ideal is
principal, which proves the converse.

For the failure of $\MU\le2$: over the wild ring
$\F_q[u,w]/\langle u^2,w^2\rangle$ the module $\mathfrak m\oplus\mathfrak
m\subset \Bq_j^2$ requires more than two generators; explicit submodules
with $\MU=3$ are constructed in Theorem~\ref{thm:nonprincipal}.
\end{proof}

\begin{remark}[the CRT step is what opens $n=5,6$]
Earlier component routines for $\Bq$ covered only three regimes: split
semisimple, $n=p^v$ with $t=1$, and ``wild'' $n=p^v$ with $t\ge2$. The general
decomposition \eqref{eq:crt} removes that restriction, which is what makes the
families $(q,t,n)=(4,1,5)$, $(4,1,6)$ and $(8,1,3)$ reachable; they are
precisely the families from which two of the headline codes in
Section~\ref{sec:cmp} come.
\end{remark}

\subsection{The reversal operator}

\begin{definition}\label{def:rev}
Let $\rev:\Rn\to\Rn$ be the $\F_q$-linear map
$\rev\bigl(\sum_ex_eX^e\bigr)=\sum_ex_{-e}X^e$, the indices being read modulo
$n$.
\end{definition}

\begin{lemma}\label{lem:rev}
$\rev$ is a ring \textbf{automorphism} of $\Rn$ (not an anti-automorphism), and
$\rev=\id$ on $\Rn$ if and only if $n\le2$.
\end{lemma}

\begin{proof}
Multiplicativity is the computation
$\rev(ab)=\sum_e\bigl(\sum_{i+j=e}a_ib_j\bigr)X^{-e}
=\sum_{i,j}a_ib_jX^{-i-j}=\rev(a)\rev(b)$; note that this uses commutativity
of $X$ with the coefficients but no commutativity of the coefficients with each
other, and it is an automorphism because the index map $e\mapsto -e$ is an
involution. If $n\le2$ then $-e\equiv e\pmod n$ for all $e$, so $\rev=\id$;
conversely, if $n\ge3$ then $\rev(X)=X^{-1}\ne X$.
\end{proof}

\subsection{Worked examples I: the ring and the CRT decomposition}

\begin{example}[the ring $R$ is not commutative, and $\theta_s$ is an
anti-automorphism]\label{ex:ring}
Take $q=4$, $t=2$, so that $A_2=\F_4[u]/\langle u^2\rangle$ and
$R=M_2(A_2)$. By Lemma~\ref{lem:basic}(b),
$|R|=q^{4t}=4^8=65536$ and $\dim_{\F_4}R=8$; the centre is $A_2\cdot I$ and
the Jacobson radical is $M_2(uA_2)$, so $R$ has exactly $t+1=3$ two-sided
ideals, namely $R$, $uR$ and $0$.

Let $E_{kl}$ denote the matrix units and take $s=0$, so that
$\theta_0(X)=X^{\mathsf T}$. Then
\[
  E_{12}E_{21}=E_{11},\qquad\text{so}\qquad
  \theta_0(E_{12}E_{21})=E_{11},
\]
whereas
\[
  \theta_0(E_{12})\,\theta_0(E_{21})=E_{21}E_{12}=E_{22}\ne E_{11},\qquad
  \theta_0(E_{21})\,\theta_0(E_{12})=E_{12}E_{21}=E_{11}.
\]
So $\theta_0(ab)=\theta_0(b)\theta_0(a)$ and $\theta_0$ is an
\emph{anti}-automorphism, not an automorphism. Over a commutative base ring
this distinction disappears entirely --- $\theta_s$ degenerates to $\sigma_s$
--- which is why the whole of Section~\ref{sec:cyclic} has no commutative
analogue.
\end{example}

\begin{example}[a wild CRT component]\label{ex:crt}
Take $(q,t,n)=(4,2,2)$. In characteristic $2$,
\[
  X^2-1=(X-1)^2\quad\text{over }\F_4,
\]
so the factorisation \eqref{eq:crt} has a single factor $g_1=X-1$ of degree
$d_1=1$ with multiplicity $e_1=2$. Hence
\[
  \Bq=A_2[X]/\langle (X-1)^2\rangle\;\cong\;\F_4[u,w]/\langle u^2,w^2\rangle,
\]
a local ring that is \textbf{not} a chain ring: its maximal ideal is
$\mathfrak m=(u,w)$, which needs two generators, and
\[
  \mathfrak m^2=(uw),\qquad \mathfrak m^3=0,\qquad
  \dim_{\F_4}\mathfrak m=3,\ \dim_{\F_4}\mathfrak m^2=1,\qquad
  \dim_{\F_4}\ann(\mathfrak m)=2 .
\]
This is the ring in which $\MU_{\Bq}(W)\le2$ fails
(Lemma~\ref{lem:components}), and it is the source of all the witnesses of
Theorem~\ref{thm:nonprincipal}. The whole family has $|\Bq|=4^4=256$ and
exactly $1221$ $\Bq$-submodules of $\Bq^2$ --- the number $1221$ in the census
of Section~\ref{sec:cert} --- each corresponding to one cyclic code of length
$2$ over $M_2(A_2)$.
\end{example}

% =====================================================================
\section{Cyclic codes and Galois duality}
\label{sec:cyclic}
% =====================================================================

\subsection{The $\theta_s$-pairing}\label{subsec:pairing}

\begin{definition}\label{def:pairing}
For $x=(x_0,\dots,x_{n-1})$, $y=(y_0,\dots,y_{n-1})\in R^n$ put
\[
  B_s(x,y)=\sum_{e=0}^{n-1}\theta_s(x_e)y_e\;\in R,
  \qquad
  \bb_s(x,y)=\Lambda\bigl(B_s(x,y)\bigr)\;\in\F_q ,
\]
and for a code $C\le\Rn$ put
\[
  C^{\perp_s}=\{y\in R^n:\bb_s(x,y)=0\ \text{for all }x\in C\}.
\]
We call $C^{\perp_s}$ the \emph{$s$-Galois dual} of $C$.
\end{definition}

\begin{lemma}[sesquilinearity]\label{lem:sesqui}
$\bb_s$ is additive in both arguments, $\F_p$-bilinear,
$\F_q$-linear in the second argument and $\sigma_s$-semilinear in the first:
\[
  \bb_s(\alpha x,y)=\sigma_s(\alpha)\,\bb_s(x,y),\qquad
  \bb_s(x,\alpha y)=\alpha\,\bb_s(x,y)\qquad(\alpha\in\F_q).
\]
\end{lemma}

\begin{proof}
$\theta_s$ is additive and $\sigma_s$-semilinear
(Lemma~\ref{lem:theta}(d)); $B_s$ is therefore additive in each argument and
satisfies $B_s(\alpha x,y)=\sigma_s(\alpha)B_s(x,y)$ (because $\alpha$ is
scalar, so it commutes with the transpose and is then moved by $\sigma_s$) and
$B_s(x,\alpha y)=B_s(x,y)\alpha$. Applying $\Lambda$, which is $\F_q$-linear,
gives the claims; note that the second identity uses
$\Lambda(z\alpha)=\alpha\Lambda(z)$, which holds because $\Lambda$ is
$\F_q$-linear.
\end{proof}

The slot convention --- $\theta_s$ acting on the \emph{first} argument --- is
immaterial for the resulting dual: placing $\theta_s$ on the second argument
produces the same subspace, because the two pairings differ by the
$\F_q$-semilinear involution $\sigma_s$ composed with transposition, which does
not change vanishing. We fix the convention of
Definition~\ref{def:pairing} throughout.

\subsection{The duality theorem}

\begin{theorem}[Galois duality on right modules]\label{thm:duality}
Let $C\le\Rn$ be a right $\Rn$-submodule. Then:
\begin{enumerate}[label=\textnormal{(\alph*)},leftmargin=*]
\item $C^{\perp_s}$ is a right $\Rn$-submodule;
\item $\dim_{\F_q}C+\dim_{\F_q}C^{\perp_s}=\Nlen$, equivalently
  $|C|\cdot|C^{\perp_s}|=|R|^n$;
\item $(C^{\perp_s})^{\perp_s}=\sigma_{2s}(C)$;
\item if $C$ is cyclic (a principal right ideal) then so is $C^{\perp_s}$.
\end{enumerate}
\end{theorem}

\begin{proof}
(a) Let $y\in C^{\perp_s}$ and $a\in\Rn$; we must show $ya\in C^{\perp_s}$.
For $x\in C$ we have $xa\in C$ (as $C$ is a right ideal), hence
\[
  0=\bb_s(xa,y)=\Lambda\Bigl(\sum_e\theta_s\bigl((xa)_e\bigr)y_e\Bigr).
\]
We claim $\bb_s(x,ya)=\bb_s(xa,y)$. Indeed, writing
$(xa)_e=\sum_{i+j=e}x_ia_j$ (indices modulo $n$),
\[
  B_s(xa,y)=\sum_e\theta_s\Bigl(\sum_{i+j=e}x_ia_j\Bigr)y_e
  =\sum_{i,j}\theta_s(a_j)\theta_s(x_i)y_{i+j},
\]
while
\[
  B_s(x,ya)=\sum_i\theta_s(x_i)(ya)_i
  =\sum_{i,j}\theta_s(x_i)y_{i+j}a_j .
\]
These are not literally equal in the non-commutative ring $R$, but their images
under $\Lambda$ agree: by Lemma~\ref{lem:lambda}, $\Lambda$ is symmetric, so
$\Lambda(\theta_s(a_j)\theta_s(x_i)y_{i+j})
=\Lambda(\theta_s(x_i)y_{i+j}\theta_s(a_j))$, and summing over $i,j$ gives
$\bb_s(xa,y)=\bb_s(x,ya)$ up to applying $\sigma_s$ to the entries of $a_j$
inside a symmetric trace --- explicitly,
\[
  \bb_s(x,ya)=\sum_{i,j}\Lambda\bigl(\theta_s(x_i)y_{i+j}a_j\bigr)
  =\sum_{i,j}\Lambda\bigl(a_j\theta_s(x_i)y_{i+j}\bigr)
  =\sum_{i,j}\sigma_s\!\bigl(\Lambda(\theta_s(a_j)\theta_s(x_i)y_{i+j})\bigr),
\]
the last step because $\Lambda(\sigma_s(z))=\sigma_s(\Lambda(z))$ and
$\Lambda$ is symmetric; since $\bb_s(xa,y)=0$ and $\sigma_s$ is injective,
$\bb_s(x,ya)=0$. Hence $ya\in C^{\perp_s}$.

(b) By Lemma~\ref{lem:sesqui} the map
\[
  \psi_C:R^n\longrightarrow C^{\vee},\qquad
  y\longmapsto\bigl(x\mapsto\bb_s(x,y)\bigr),
\]
where $C^{\vee}$ denotes the $\F_q$-dual \emph{with the $\sigma_s$-semilinear
action in the first slot}, is $\F_q$-linear and surjective
(non-degeneracy of $\Lambda$, Lemma~\ref{lem:lambda}, applied coordinate-wise).
Its kernel is $C^{\perp_s}$, so
$\dim_{\F_q}C^{\perp_s}=\Nlen-\dim_{\F_q}C$.

(c) By (b) applied twice,
$\dim(C^{\perp_s})^{\perp_s}=\Nlen-\dim C^{\perp_s}=\dim C=\dim\sigma_{2s}(C)$,
so it suffices to prove one inclusion. For $x\in C$ and $y\in C^{\perp_s}$ we
have $\bb_s(y,\sigma_{2s}(x))=0$: indeed
$\bb_s(y,\sigma_{2s}(x))=\sigma_s(\bb_s(x,y))$ by Lemma~\ref{lem:theta}(c) and
Lemma~\ref{lem:sesqui}, and $\bb_s(x,y)=0$. Hence
$\sigma_{2s}(C)\subseteq(C^{\perp_s})^{\perp_s}$ and dimensions agree.

(d) If $C=g\Rn$ then $C^{\perp_s}$ is a right ideal by (a); its cyclicity
follows from the explicit description of Theorem~\ref{thm:dualformula} below,
which exhibits it as the annihilator of a single element.
\end{proof}

\begin{remark}
Theorem~\ref{thm:duality}(a) is the statement that fails for the plain product
on right modules (Lemma~\ref{lem:failure}): the transpose in $\theta_s$ is
exactly what converts the left action into a right action.
\end{remark}

\subsection{Explicit duals: annihilators and the reversal}

For a subset $I\subseteq\Rn$ write
$\ann_r(I)=\{a:xa=0\ \forall x\in I\}$ and
$\ann_l(I)=\{a:ax=0\ \forall x\in I\}$.

\begin{lemma}\label{lem:ann}
$\ann_r(\sum_i\Rn x_i)=\bigcap_i\ann_r(x_i)$; for a right ideal
$I$, $\ann_r(I)$ is determined by any generating set, whereas $\ann_l(I)$ is
not. Moreover $\theta_s(\ann_r(I))=\ann_l(\theta_s(I))$.
\end{lemma}

\begin{proof}
The first identity is the definition of an annihilator of a sum. The asymmetry
between left and right is a right-ideal phenomenon: a right ideal is generated
as a right ideal, and $\ann_r$ sees right multiplication, which is the side on
which the generators act. The last identity: $a\in\ann_r(I)$ iff $xa=0$ for all
$x\in I$; applying the anti-automorphism $\theta_s$ gives
$\theta_s(a)\theta_s(x)=0$ for all $x\in I$, i.e.
$\theta_s(a)\in\ann_l(\theta_s(I))$; conversely by applying $\theta_s^{-1}$.
\end{proof}

\begin{theorem}[dual formula]\label{thm:dualformula}
Let $C\le\Rn$ be a right $\Rn$-submodule. Then
\[
  C^{\perp_s}=\rev\Bigl(\ann_r\bigl(\theta_s(C)\bigr)\Bigr).
\]
In particular, if $C=g\Rn$ is principal then
$C^{\perp_s}=\rev(\ann_r(\theta_s(g)))$ and $C^{\perp_s}$ is cyclic.
\end{theorem}

\begin{proof}
Write $x=\sum_ex_eX^e$, $y=\sum_ey_eX^e$. Then
\[
  B_s(x,y)=\sum_e\theta_s(x_e)y_e .
\]
On the other hand, the constant coefficient of the product
$\theta_s(x)\cdot y$ (computed in $\Rn$) is
\[
  \bigl[\theta_s(x)y\bigr]_0
  =\sum_{i+j\equiv0}\theta_s(x_i)y_j
  =\sum_i\theta_s(x_i)y_{-i}
  =\sum_e\theta_s(x_e)\,\bigl(\rev y\bigr)_e ,
\]
because $(\rev y)_i=y_{-i}$. Hence
\[
  \bb_s(x,y)=\Lambda\Bigl(\bigl[\theta_s(x)\,\rev y\bigr]_0\Bigr).
\]
Now $y\in C^{\perp_s}$ iff $\bb_s(x,y)=0$ for all $x\in C$, iff
$\Lambda([\theta_s(x)v]_0)=0$ for all $x\in C$, where $v=\rev y$. Since the
constant coefficient of $zv$ is $\F_q$-linear in $z$ and $\Lambda$ is
non-degenerate, and since the set of $z$ arising is $\theta_s(C)$, this is
equivalent to $\Lambda([zv]_0)=0$ for all $z\in\theta_s(C)$.

It remains to identify this condition with $v\in\ann_r(\theta_s(C))$. Write
$z=\sum_iz_iX^i$, $v=\sum_jv_jX^j$. Then
$[zv]_0=\sum_{i+j\equiv0}z_iv_j=\sum_iz_iv_{-i}$, and, using the symmetry of
$\Lambda$,
\[
  \Lambda([zv]_0)=\sum_i\Lambda(z_iv_{-i})
  =\sum_i\Lambda(v_{-i}z_i).
\]
This vanishes for all $z\in\theta_s(C)$ exactly when $v$ right-annihilates
$\theta_s(C)$: the coefficient extraction ``take the constant coefficient and
apply $\Lambda$'' is the trace pairing that detects annihilation, and the index
reversal is precisely $\rev$. Therefore $v\in\ann_r(\theta_s(C))$ and
$y=\rev(v)$, which is the assertion.

For the principal case, $\ann_r(\theta_s(C))=\ann_r(\theta_s(g)\Rn)
=\ann_r(\theta_s(g))$ by Lemma~\ref{lem:ann}, and an annihilator of a single
element is a right ideal generated by one element in this ring, so
$C^{\perp_s}$ is cyclic.
\end{proof}

\begin{corollary}[the reversal is load-bearing exactly for $n\ge3$]
\label{cor:rev}
For $n\le2$ the operator $\rev$ is the identity, so
$C^{\perp_s}=\ann_r(\theta_s(C))$; for $n\ge3$ dropping $\rev$ changes the dual.
\end{corollary}

\begin{proof}
Immediate from Lemma~\ref{lem:rev} and Theorem~\ref{thm:dualformula}.
\end{proof}

\subsection{Equality of the two orthogonals}

There are two natural duals of $C\le R^n$: the $R$-valued one,
$C^{\perp,R}=\{y:B_s(x,y)=0\ \forall x\in C\}$, and the $\F_q$-valued one
$C^{\perp_s}$ of Definition~\ref{def:pairing}.

\begin{proposition}\label{prop:twoorth}
$C^{\perp,R}=C^{\perp_s}$ for every $C$ and every $s$.
\end{proposition}

\begin{proof}
The inclusion $C^{\perp,R}\subseteq C^{\perp_s}$ is immediate since
$\bb_s=\Lambda\circ B_s$. For the converse, both are $\F_q$-subspaces of $R^n$;
$\dim_{\F_q}C^{\perp_s}=\Nlen-\dim C$ by Theorem~\ref{thm:duality}(b), and the
same dimension count applies to $C^{\perp,R}$: the map
$y\mapsto(x\mapsto B_s(x,y))$ into $\Hom_{\F_q}(C,R)$ is $\F_q$-linear with
kernel $C^{\perp,R}$, and it is surjective because $\Lambda$ is
non-degenerate and $B_s$ is non-degenerate in $y$. Hence the two subspaces have
the same dimension and the inclusion is an equality.
\end{proof}

\subsection{Three regimes}\label{subsec:regimes}

The parameter $s$ produces three genuinely different situations.

\begin{definition}\label{def:regimes}
Let $q=p^r$.
\begin{enumerate}[label=\textnormal{(\arabic*)},leftmargin=*]
\item The \emph{general Galois regime}: no condition on $s$.
\item The \emph{reflexive regime}: $r\mid 2s$, i.e. $\sigma_{2s}=\id$.
\item The \emph{per-code invariant regime}: the individual code $C$ satisfies
  $\sigma_{2s}(C)=C$.
\end{enumerate}
\end{definition}

\begin{proposition}\label{prop:regimes}
$(2)\Rightarrow(3)$, and the converse fails. In the reflexive regime
$(C^{\perp_s})^{\perp_s}=C$ as subspaces; in general only
$(C^{\perp_s})^{\perp_s}=\sigma_{2s}(C)$.
\end{proposition}

\begin{proof}
The implication is clear. For the failure of the converse, take
$q=8$ ($r=3$), $s=1$, so $2s=2$ and $r\nmid2$: at $(q,t,n)=(8,1,2)$ exactly
$14$ of the $92$ codes satisfy $\sigma_2(C)=C$, a strict subset. The two
displayed identities are Theorem~\ref{thm:duality}(c) together with
$\sigma_{2s}=\id$.
\end{proof}

\begin{remark}
The three regimes matter because the symmetry of the orthogonality
\emph{relation} holds exactly in the reflexive regime, whereas --- as we show in
Theorem~\ref{thm:hullso} --- the self-orthogonality of the hull holds in
\emph{all} of them. This is the point at which the present treatment departs
from treatments that assume reflexivity.
\end{remark}

\subsection{Worked examples II: duality in a reflexive and in a
non-reflexive regime}

\begin{example}[a Galois dual computed both ways, at $(q,t,n)=(4,1,1)$]
\label{ex:dual}
Here $A_1=\F_4$, $R=M_2(\F_4)$, $\Bq=\F_4[X]/\langle X-1\rangle\cong\F_4$, and
$\Nlen=4tn=4$. By \eqref{eq:morita}, right ideals of $\Rn=M_2(\F_4)$
correspond to $\F_4$-subspaces of $\F_4^2$; there are $7$ of them, which is
the entry $7$ in the census table of Section~\ref{sec:cert}.

Take $W=\F_4\cdot(1,0)^{\mathsf T}\le\F_4^2$. The corresponding right ideal is
\[
  C=I_W=\{M\in M_2(\F_4):\text{both columns of }M\text{ lie in }W\}
  =\left\{\begin{pmatrix}a&b\\0&0\end{pmatrix}:a,b\in\F_4\right\},
\]
of $\F_4$-dimension $k=2$. Take $s=0$; then $\bb_0(X,Y)=\tr(X^{\mathsf T}Y)
=\sum_{k,l}X_{kl}Y_{kl}$, i.e. the transported form is the componentwise dot
product on $\F_4^2$, and $W^{\perp_0}=\F_4\cdot(0,1)^{\mathsf T}$. Hence
\[
  C^{\perp_0}=I_{W^{\perp_0}}
  =\left\{\begin{pmatrix}0&0\\a&b\end{pmatrix}:a,b\in\F_4\right\},
\]
again of dimension $2$. Three checks:
\begin{enumerate}[label=\textnormal{(\roman*)},leftmargin=*]
\item $C\cap C^{\perp_0}=\{0\}$, so $C$ is LCD (Proposition~\ref{prop:lcdso});
\item $\dim C+\dim C^{\perp_0}=2+2=4=\Nlen$, as
  Theorem~\ref{thm:duality}(b) requires --- and $|C|\cdot|C^{\perp_0}|
  =16\cdot16=256=|\Rn|$, which is exactly what fails for the plain product
  (Lemma~\ref{lem:failure});
\item $C\oplus C^{\perp_0}=\Rn$, so this is also the LCD direct-sum
  splitting.
\end{enumerate}
Every step above is a $2\times2$ matrix computation over $\F_4$ and can be
checked by hand; it is the smallest instance of the whole dictionary.
\end{example}

\begin{example}[a non-reflexive regime, where $(C^{\perp_s})^{\perp_s}\ne C$]
\label{ex:nonreflexive}
Take $q=8$, so $r=3$ and $\sigma_s(a)=a^{2^s}$ has order
$3/\gcd(3,s)$. Take $s=1$. Then $2s=2$ and $r\nmid 2s$, so we are in the
general Galois regime and not the reflexive one:
$\sigma_2=\sigma_1^2\ne\id$. Consequently
\[
  (C^{\perp_1})^{\perp_1}=\sigma_2(C)\ne C
\]
for any code $C$ not invariant under $\sigma_2$. At $(q,t,n)=(8,1,2)$ there
are $92$ nonzero cyclic codes with invariants computed, of which exactly $14$
satisfy $\sigma_2(C)=C$; for those $14$ the per-code invariant regime holds
even though the reflexive regime does not.

This example is the reason Theorem~\ref{thm:hullso} is proved by a
one-sided argument: in this regime the orthogonality \emph{relation} is not
symmetric, so any argument that swaps the two slots of $\bb_s$ would be
invalid here. For $q=4$ and $q=9$ one has $r=2$, so $r\mid 2s$ for every $s$
and every regime is reflexive; the non-reflexive case arises first at $q=8$.
\end{example}

% =====================================================================
\section{Galois hull structure and main algebraic results}
\label{sec:hulls}
% =====================================================================

\begin{definition}\label{def:hull}
For $C\le\Rn$ and $s\ge0$, the \emph{$s$-Galois hull} of $C$ is
\[
  \Hull_s(C)=C\cap C^{\perp_s}.
\]
We write $h_s(C)=\dim_{\F_q}\Hull_s(C)$.
\end{definition}

\subsection{Structure of the hull}

\begin{theorem}[hull structure]\label{thm:hullstruct}
Let $C\le\Rn$ be a right $\Rn$-submodule. Then:
\begin{enumerate}[label=\textnormal{(\alph*)},leftmargin=*]
\item $\Hull_s(C)$ is a right $\Rn$-submodule;
\item $\Hull_s(C)$ is invariant under the cyclic shift $T$, i.e.
  $T\bigl(\Hull_s(C)\bigr)=\Hull_s(C)$; it is a right $\Rn$-submodule but
  \textbf{not} in general \emph{principal} --- that is
  Theorem~\ref{thm:nonprincipal} below;
\item $\Hull_s(C)$ is \textbf{not} in general a left ideal, hence not a
  two-sided ideal;
\item $\Hull_s(\sigma_kC)=\sigma_k\Hull_s(C)$ for every $k$, and
  $\Hull_s(\sigma_{2s}C)=\sigma_{2s}\Hull_s(C)$;
\item in the reflexive regime ($r\mid 2s$), $\Hull_s(C^{\perp_s})
  =\Hull_s(C)$ as subspaces; in general
  $\Hull_s(C^{\perp_s})=C^{\perp_s}\cap\sigma_{2s}(C)$.
\end{enumerate}
\end{theorem}

\begin{proof}
(a) Intersection of two right $\Rn$-submodules
(Theorem~\ref{thm:duality}(a)); note that no principality hypothesis is used.

(b) $T$-invariance. Since $T$ and $\theta_s$ both act entrywise on the
coefficient tuple (the one by a cyclic index shift, the other by transpose and
Frobenius), they commute: $T\theta_s=\theta_sT$. The pairing
$\bb_s(x,y)=\Lambda\bigl(\sum_e\theta_s(x_e)y_e\bigr)$ is a sum over
positions, so $\bb_s(Tx,Ty)=\bb_s(x,y)$. Now $C$ is $T$-invariant and $T$ is a
bijection, so for $y\in R^n$,
\[
  y\in C^{\perp_s}
  \iff \bb_s(x,y)=0\ \forall x\in C
  \iff \bb_s(Tx,Ty)=0\ \forall x\in C
  \iff \bb_s(x',Ty)=0\ \forall x'\in C
  \iff Ty\in C^{\perp_s}.
\]
Hence $C^{\perp_s}$ is $T$-invariant, and $T$ preserves the intersection
$C\cap C^{\perp_s}$. The \emph{principality} of the hull is not asserted here
and is false in general (Theorem~\ref{thm:nonprincipal}); what (a) and (b)
give is that $\Hull_s(C)$ is a right $\Rn$-submodule, which is all that
Section~\ref{sec:gray} and Section~\ref{sec:quantum} use.

(c) An explicit witness: at $(q,t,n)=(4,1,1)$, $s=0$, with
$g=\begin{pmatrix}0&1\\0&1\end{pmatrix}$, the hull $\Hull_0(g\Rn)$ is a right
ideal that is not closed under left multiplication by $R$.

(d) $\sigma_k$ is a ring automorphism fixing $X$, so it preserves right ideals
and commutes with $\theta_s$ up to the index shift
$\sigma_k\theta_s=\theta_s\sigma_k$; hence it carries $C^{\perp_s}$ to
$(\sigma_kC)^{\perp_s}$ and therefore the intersection to the intersection. The
second identity is the case $k$ replaced by $2s$, using
Theorem~\ref{thm:duality}(c).

(e) By Theorem~\ref{thm:duality}(c),
$(C^{\perp_s})^{\perp_s}=\sigma_{2s}(C)$, so
$\Hull_s(C^{\perp_s})=C^{\perp_s}\cap\sigma_{2s}(C)$. In the reflexive regime
$\sigma_{2s}=\id$, giving $\Hull_s(C^{\perp_s})=C^{\perp_s}\cap C=\Hull_s(C)$.
\end{proof}

\begin{proposition}[Morita hull formula and the principality criterion]
\label{prop:moritahull}
Under \eqref{eq:morita}, let $W=W(C)\le\Bq^2$ be the $\Bq$-submodule
corresponding to $C$ and $W'=W(C^{\perp_s})$. Then
\[
  \Hull_s(C)=I_{W\cap W'},
\]
the right ideal of $\Rn$ corresponding to $W\cap W'$, and
$\Hull_s(C)$ is principal if and only if $\MU_{\Bq}(W\cap W')\le2$.
\end{proposition}

\begin{proof}
Under the Morita correspondence, intersection of right ideals corresponds to
intersection of the associated submodules of $\Bq^2$, which gives the displayed
formula. The principality criterion is the Morita dictionary applied to the
minimal number of generators: a right ideal $I_U$ is principal iff $U$ is
generated by at most $2$ elements as a $\Bq$-module (the two standard basis
vectors of $\Bq^2$ being the generators of $\Rn$ as a right module).
\end{proof}

\subsection{Non-principal hulls: the conjecture O-3.1 is false}

Over a finite chain ring every submodule of $S^2$ is free of rank at most $2$,
hence $2$-generated. Combined with Lemma~\ref{lem:components} this suggests the
following natural conjecture, which was the open problem O-3.1 of this project.

\begin{quote}\textbf{Conjecture O-3.1 (false, as stated).} For every cyclic
code $C$ over $R=M_2(A_t)$ and every $s$, the Galois hull $\Hull_s(C)$ is a
principal right ideal of $\Rn$.
\end{quote}

Two readings must be separated, because only the first is refuted here:
\begin{description}
\item[(O-3.1a)] for \emph{every} cyclic code $C$ (i.e. every right ideal of
  $\Rn$, principal or not) the hull is principal. --- \textbf{Refuted},
  Theorem~\ref{thm:nonprincipal}.
\item[(O-3.1b)] for every \emph{principal} $C=g\Rn$ the hull is principal.
  --- \textbf{Open}; see Remark~\ref{rem:o31b}.
\end{description}

\begin{theorem}\label{thm:nonprincipal}
Conjecture O-3.1 is false. Moreover:
\begin{enumerate}[label=\textnormal{(\alph*)},leftmargin=*]
\item $\MU_{\Bq}(\Hull_s(C))\ge3$ requires $t\ge2$ \textbf{and} $p\mid n$;
\item the bound $\MU_{\Bq}(W)\le2$ is false over a non-chain local
  component, and it fails for a hull.
\end{enumerate}
\end{theorem}

\begin{proof}
(a) By Lemma~\ref{lem:components}, $\MU_{\Bq}(W)=\max_j\MU_{\Bq_j}(W_j)$, and
$\Bq_j$ is a chain ring iff $t=1$ or $e_j=1$. If $t=1$ then every component is a
chain ring and every $W_j$ is $2$-generated, so $\MU_{\Bq}(W)\le2$. If
$p\nmid n$ then $X^n-1$ is separable over $\F_q$, so every $e_j=1$ and again
every component is a chain ring. Hence $\MU_{\Bq}(W)\ge3$ forces $t\ge2$ and
$p\mid n$ (equivalently: some $e_j\ge2$, which happens exactly when $p\mid n$,
since $X^n-1=(X^{n'}-1)^{p^v}$ with $n=p^vn'$ and $p\nmid n'$).

(b) Let $\Bq_j\cong\F_q[u,w]/\langle u^2,w^2\rangle$ be a wild component. Its
maximal ideal is $\mathfrak m=(u,w)$, $\mathfrak m^2=(uw)$,
$\mathfrak m^3=0$, and $\ann_{\Bq_j^2}(\mathfrak m)$ has $\F_q$-dimension $2$.
Consider the family of submodules
\begin{equation}\label{eq:witnessfamily}
  W=\mathfrak m^ie_1\oplus U,\qquad U\le\ann_{\Bq^2}(\mathfrak m),
\end{equation}
where $e_1$ is the first standard basis vector of $\Bq^2$. Then
$\mathfrak mW=\mathfrak m^{i+1}e_1$, so by Nakayama
\[
  \MU_{\Bq}(W)=\dim_{\F_q}\mathfrak m^i-\dim_{\F_q}\mathfrak m^{i+1}
  +\dim_{\F_q}U .
\]
Taking $i=0$ (so $\dim\mathfrak m^0-\dim\mathfrak m^1=\dim\Bq-\dim\mathfrak m
=4-3=1$ over $\F_q$ for $t=2$, $e_j=2$) and $\dim U=2$ gives
$\MU_{\Bq}(W)=3$. Enumerating the deterministic family \eqref{eq:witnessfamily}
exhaustively in the wild families produces submodules whose Galois hull has
$\MU_{\Bq}=3$; see Section~\ref{subsec:witnesses} for the explicit witnesses.
\end{proof}

\begin{remark}
The family \eqref{eq:witnessfamily} is finite and completely deterministic, so
the witnesses are reproducible without search. The use of $e_1$ and of
$\ann(\mathfrak m)$ is essential: the attempted construction
$W=\mathfrak m^ie_1\oplus U$ with $U$ arbitrary does \emph{not} give a
$\Bq$-submodule, and the direct ``complement-lift'' construction fails.
\end{remark}

\subsection{The witnesses}\label{subsec:witnesses}

The following families have been examined; in each, the deterministic family
\eqref{eq:witnessfamily} was enumerated completely and the Galois hull
$\Hull_s(W)$ computed by the Morita route, with $\MU_{\Bq}$ of the hull
recomputed directly (not inferred).

\begin{table}[h]
\centering
\caption{Wild families: maximal ideal structure and the best witness found.
Here $\dim\mathfrak m^i$ lists $\dim_{\F_q}\mathfrak m^i$ for $i\ge1$.}
\label{tab:witnesses}
\begin{tabular}{@{}ccccccc@{}}
\toprule
$(q,t,n,s)$ & $\Nlen$ & $\dim\mathfrak m,\dim\mathfrak m^2,\dots$
 & $\dim\ann(\mathfrak m)$ & $\#\{W:\MU_{\Bq}(W)\ge3\}$
 & $\MU_{\Bq}(\Hull_s)$ & $\Hull_s=W$? \\
\midrule
$(4,2,2,0)$ & 16 & $3,1$ & 2 & 12 & 2 & --- \\
$(4,2,2,1)$ & 16 & $3,1$ & 2 & 12 & 2 & --- \\
$(8,2,2,0)$ & 16 & $3,1$ & 2 & 56 & 2 & --- \\
$(8,2,2,1)$ & 16 & $3,1$ & 2 & 56 & 2 & --- \\
$(9,2,3,0)$ & 24 & $5,3,1$ & 2 & 144 & \textbf{3} & yes \\
$(9,2,3,1)$ & 24 & $5,3,1$ & 2 & 144 & \textbf{3} & yes \\
$(4,2,4,0)$ & 32 & $7,5,3,1$ & 2 & 36 & \textbf{3} & no \\
$(4,2,4,1)$ & 32 & $7,5,3,1$ & 2 & 36 & \textbf{3} & no \\
\bottomrule
\end{tabular}
\end{table}

In addition, an exhaustive enumeration of \emph{every} code of \emph{every}
enumerable family (families of sizes $7,33,33,1221,343,144,93,35937$) found
four hulls with $\MU_{\Bq}=3$ at $(4,2,2)$ --- one at $s=0$ and three at $s=1$
--- which is the original refutation of O-3.1.

\begin{remark}[the sharper question (O-3.1b) remains open]\label{rem:o31b}
Every witness found --- the four at $(4,2,2)$ and the four at $(9,2,3)$ and
$(4,2,4)$ --- has $\MU_{\Bq}(W)=3$ for the \emph{code} as well as for the
hull. Equivalently, every known witness is a code that is itself
non-principal. The four $(4,2,2)$ witnesses and the two $(9,2,3)$ witnesses
additionally satisfy $\Hull_s(W)=W$, so the hull is non-principal because the
code is.

It is therefore \textbf{open} whether a \emph{principal} cyclic code can have a
non-principal Galois hull in the wild regime. We do not claim it either way,
and we note that the two questions are frequently conflated. What
Theorem~\ref{thm:nonprincipal} settles is that hull principality cannot be
assumed, so no downstream argument may use it.
\end{remark}

\begin{remark}[a disambiguation]\label{rem:mub}
The number $\MU_B$ is attached to a \emph{module}. The value $\MU_B=4$ reported
in earlier work of this project refers to $\MU_B(W)$ for the \emph{code}
$W$ at $(4,2,2,1)$; the value constrained by the hull census is
$\MU_B(\Hull_s(W))$, whose observed range is $\{0,1,2,3\}$ and whose maximum
over the extended enumeration is $3$. The two numbers are not in conflict.
\end{remark}

\subsection{What does not exist}

\begin{theorem}[no closed hull formula]\label{thm:noclosed}
There is no universal closed formula of ``gcd type'' for the Galois hull of a
cyclic code over $R=M_2(A_t)$: that is, there is no expression in the
generators of $C$, built from the ring operations and a finite fixed set of
polynomial operations, that yields a generator of $\Hull_s(C)$ for all $C$ and
all $s$.
\end{theorem}

\begin{proof}
If such a formula existed, every hull would be principal, since the output of
the formula is (by construction) a single element generating the hull. This
contradicts Theorem~\ref{thm:nonprincipal}.
\end{proof}

\begin{remark}
One should not read Theorem~\ref{thm:noclosed} as a negative result about the
\emph{usefulness} of hulls. The Morita formula of
Proposition~\ref{prop:moritahull} is a perfectly effective way to compute
$\Hull_s(C)$: it is a submodule intersection in $\Bq^2$, and via the CRT
decomposition \eqref{eq:crt} it is computed componentwise. What does not exist
is a generator-polynomial formula of the familiar commutative kind.
\end{remark}

\subsection{LCD, self-orthogonal and self-dual codes}

\begin{proposition}\label{prop:lcdso}
Let $C\le\Rn$ with Morita module $W=W(C)$ and $W'=W(C^{\perp_s})$. Then
\begin{enumerate}[label=\textnormal{(\alph*)},leftmargin=*]
\item $C$ is LCD ($\Hull_s(C)=\{0\}$) iff $W\cap W'=\{0\}$ iff
  $C\oplus C^{\perp_s}=R^n$;
\item $C\subseteq C^{\perp_s}$ iff $W\subseteq W'$; for $C=g\Rn$ principal
  this is equivalent to $\theta_s(\rev g)\,g=0$;
\item $C=C^{\perp_s}$ iff $C\subseteq C^{\perp_s}$ and $2\dim_{\F_q}C=\Nlen$;
  for $C=g\Rn$ this is equivalent to $\im g=\ker\theta_s(\rev g)$.
\end{enumerate}
\end{proposition}

\begin{proof}
Immediate from Proposition~\ref{prop:moritahull} and the dimension count of
Theorem~\ref{thm:duality}(b); the principal-case criteria follow by translating
$W=\im g$ and $W'=\ker\theta_s(\rev g)$ through
Theorem~\ref{thm:dualformula}.
\end{proof}

\begin{remark}[two conjectures, stated as such]
Let $D=C^{\perp_s}$. It is natural to ask whether
$\Hull_s(C)\subseteq\sigma_{2s}(C)$ always holds, equivalently whether
$\Hull_s(D)=\Hull_s(D^{\perp_s})$; and whether
$\dim\Hull_s(C)=\dim\Hull_s(C^{\perp_s})$ always. Both hold in every instance
computed, and the first would follow from $\Hull_s(C^{\perp_s})\subseteq
\sigma_{2s}(C^{\perp_s})$. We state them as conjectures: neither is used
anywhere in this paper, and neither is claimed as a theorem.
\end{remark}

\subsection{Worked examples III: hulls}

\begin{example}[an LCD code, by the Morita criterion]\label{ex:lcd}
With the notation of Example~\ref{ex:dual} --- $(q,t,n)=(4,1,1)$, $s=0$ --- the
code $C=I_W$ with $W=\F_4(1,0)^{\mathsf T}$ has $W\cap W^{\perp_0}=\{0\}$, so
by Proposition~\ref{prop:moritahull}
\[
  \Hull_0(C)=I_{\{0\}}=\{0\},\qquad h_0(C)=0,
\]
and $C\oplus C^{\perp_0}=\Rn$; $C$ is LCD. The other extreme is $W=\F_4^2$, for
which $C=\Rn$ and $C^{\perp_0}=\{0\}$, again LCD.

The complete picture in this family is small enough to state: $\F_4^2$ has
$1+5+1=7$ subspaces, and a $1$-dimensional $W=\F_4v$ is self-orthogonal iff
$v\cdot v=0$. Since $v\cdot v=v_0^2+v_1^2=(v_0+v_1)^2$ in characteristic $2$,
the self-orthogonal lines are exactly those with $v_0+v_1=0$, of which there is
one, namely $v=(1,1)$. So of the $7$ codes, $6$ are LCD and one is self-dual.
\end{example}

\begin{example}[a self-orthogonal and a self-dual code]\label{ex:sd}
Take $W=\F_4\cdot(1,\omega)^{\mathsf T}$ with $\omega\in\F_4$ a root of
$\omega^2+\omega+1=0$, again at $(q,t,n)=(4,1,1)$ and $s=0$. Then
$v\cdot v=1+\omega^2=1+(\omega+1)=\omega\ne0$, so $W\nsubseteq W^{\perp_0}$
and $C=I_W$ is LCD.

Taking instead $W=\F_4\cdot(1,1)^{\mathsf T}$ gives $v\cdot v=1+1=0$, so
$W\subseteq W^{\perp_0}$. Both spaces have dimension $1$, hence
$W=W^{\perp_0}$, and therefore
\[
  C=C^{\perp_0},\qquad \Hull_0(C)=C,\qquad h_0(C)=\dim C=2=\Nlen/2 :
\]
the code is self-dual. By Proposition~\ref{prop:lcdso}(c), for a principal
$C=g\Rn$ the criterion is $\im g=\ker\theta_s(\rev g)$, which in this family
(where $\rev=\id$ by Corollary~\ref{cor:rev}) is a $4\times4$ rank
computation over $\F_4$.
\end{example}

\begin{example}[the $(9,2,3)$ witness, in full]\label{ex:witnessA}
Let $q=9$, $t=2$, $n=3$. Then $X^3-1=(X-1)^3$ over $\F_9$, so
\[
  \Bq\cong\F_9[u,w]/\langle u^2,w^3\rangle ,
\]
local with maximal ideal $\mathfrak m=(u,w)$ satisfying
\[
  \dim_{\F_9}\mathfrak m=5,\quad \dim_{\F_9}\mathfrak m^2=3,\quad
  \dim_{\F_9}\mathfrak m^3=1,\qquad \dim_{\F_9}\ann_{\Bq^2}(\mathfrak m)=2 .
\]
This is a wild component ($t\ge2$ and $e_1=3\ge2$), so by
Lemma~\ref{lem:components} the bound $\MU_{\Bq}\le2$ is not available here.
Enumerate the deterministic family \eqref{eq:witnessfamily},
\[
  W=\mathfrak m^ie_1\oplus U,\qquad U\le\ann_{\Bq^2}(\mathfrak m),
\]
completely: $144$ of its members have $\MU_{\Bq}(W)\ge3$. For each, form the
Galois hull $\Hull_s(W)$ by the Morita route and recompute
$\MU_{\Bq}(\Hull_s(W))$ directly. The outcome for $s=0$ and for $s=1$ is
\[
  \MU_{\Bq}(W)=3,\qquad \MU_{\Bq}\bigl(\Hull_s(W)\bigr)=3,\qquad
  \Hull_s(W)=W .
\]
The last equality says the code equals its own Galois hull. The corresponding
image code $D=\varphi_3(C)\subseteq\F_9^{24}$ has $k=\dim D=10$ and Euclidean
hull $h_E=8$, so route R3 gives
\[
  [[24,\,24-10-8,\,2\,;\,10-8]]_9=[[24,6,2;2]]_9 ,
\]
while route R1 gives $[[24,14,2]]_9$ with an abelian stabilizer group. The
parameters are weak; the point is that the hull is \emph{not principal}, and
yet every step of the construction applies to it unchanged. Together with
Example~\ref{ex:witnessB} this is the refutation of O-3.1 outside the
originally known family $(4,2,2)$.
\end{example}

\begin{example}[the $(4,2,4)$ witness]\label{ex:witnessB}
Let $q=4$, $t=2$, $n=4$. Then $X^4-1=(X-1)^4$ over $\F_4$ and
\[
  \Bq\cong\F_4[u,w]/\langle u^2,w^4\rangle,\qquad
  \dim_{\F_4}\mathfrak m=7,\ \dim\mathfrak m^2=5,\ \dim\mathfrak m^3=3,\
  \dim\mathfrak m^4=1,\ \ \dim\ann_{\Bq^2}(\mathfrak m)=2 .
\]
The same construction \eqref{eq:witnessfamily} has $36$ members with
$\MU_{\Bq}(W)\ge3$, and for $s=0$ and $s=1$ the best witness satisfies
\[
  \MU_{\Bq}(W)=3,\qquad \MU_{\Bq}\bigl(\Hull_s(W)\bigr)=3,\qquad
  \Hull_s(W)\ne W .
\]
Here $\Nlen=4tn=32$, $k=14$, $h_E=10$, so route R3 gives
\[
  [[32,\,32-14-10,\,2\,;\,14-10]]_4=[[32,8,2;4]]_4 .
\]
Unlike Example~\ref{ex:witnessA}, the hull is a proper superobject of the
code one started from. It is nonetheless carried unchanged through the Gray
map by Theorem~\ref{thm:grayhull}, which is precisely the statement that needs
no principality.
\end{example}

% =====================================================================
\section{Gray-map compatibility and finite-field representation}
\label{sec:gray}
% =====================================================================

This section is the bridge from the ring $R$ to $\F_q$. The map is the
coefficient expansion; the content is the identification of the form it
transports.

\subsection{The coefficient expansion}

\begin{definition}[coefficient expansion]\label{def:gray}
For $a=\sum_{a'=0}^{t-1}c_{a'}u^{a'}\in A_t$ put
$\varphi_A(a)=(c_0,\dots,c_{t-1})\in\F_q^{\,t}$. For
$X=(X_{kl})\in R=M_2(A_t)$ put
$\varphi_R(X)=(\varphi_A(X_{00}),\varphi_A(X_{01}),\varphi_A(X_{10}),
\varphi_A(X_{11}))\in\F_q^{\,4t}$, and for $x=(x_0,\dots,x_{n-1})\in R^n$,
\[
  \varphi_n(x)=\bigl(\varphi_R(x_0),\dots,\varphi_R(x_{n-1})\bigr)
  \in\F_q^{\,\Nlen},\qquad \Nlen=4tn .
\]
The flat index is $\varphi_n(x)_{(e\cdot4+k)t+a}=
\text{the $u^a$-coefficient of the $k$-th entry of }x_e$.
\end{definition}

\begin{lemma}\label{lem:philin}
$\varphi_n:R^n\to\F_q^{\Nlen}$ is an $\F_q$-linear bijection; it is additive
and satisfies $\varphi_nT=\widehat T\varphi_n$, where $T$ is the cyclic shift
on $R^n$ and $\widehat T$ is the rotation by $4t$ on $\F_q^{\Nlen}$. It is
\textbf{not} a ring homomorphism: $R$ is not commutative, and
$\varphi_R(XY)\ne\varphi_R(YX)$ in general.
\end{lemma}

\begin{proof}
$\varphi_A$ is a bijection $A_t\to\F_q^t$ by the uniqueness of the normal form;
$\varphi_R$ and $\varphi_n$ are products of it, hence bijections, and they are
$\F_q$-linear because taking the $u^a$-coefficient is. For the shift:
$T$ sends $x_e$ to $x_{e-1}$, and $\widehat T$ rotates the flat index by $4t$,
which is exactly one $R$-block; the two agree. Non-multiplicativity is the
explicit computation $\varphi_R(XY)\ne\varphi_R(YX)$ on matrix units.
\end{proof}

\subsection{The exact induced form}\label{subsec:induced}

Because $\varphi_n$ is an $\F_q$-linear bijection, there is exactly one form on
$\F_q^{\Nlen}$ for which $\varphi_n$ is a form-isometry, namely the transport of
$\bb_s$. The content of the following theorem is to identify it.

\begin{theorem}[Gram matrix of the transported form]\label{thm:gram}
Transport $\bb_s$ along $\varphi_n$:
$\langle u,v\rangle^{\widehat{}}_s:=\bb_s(\varphi_n^{-1}u,\varphi_n^{-1}v)$.
Then, writing $u^{[s]}=(\sigma_s(u_i))_i$,
\[
  \langle u,v\rangle^{\widehat{}}_s=u^{[s]}\,J\,v^{\mathsf T},
  \qquad
  J=I_{4n}\otimes H_t,\qquad
  (H_t)_{ab}=\begin{cases}1,&a+b=t-1,\\0,&\text{else.}\end{cases}
\]
The matrix $J$ is non-singular, symmetric and Hankel, $J^2=I$,
$J=I$ iff $t=1$, and $J$ is alternating iff $t$ is even.
\end{theorem}

\begin{proof}
First, $\tr(\theta_s(X)Y)=\sum_{k,l}\sigma_s(X_{kl})Y_{kl}$: indeed
$\theta_s(X)=\sigma_s(X^{\mathsf T})$ has $(k,l)$-entry
$\sigma_s(X_{lk})$, so
$\tr(\theta_s(X)Y)=\sum_{k,l}\sigma_s(X_{lk})Y_{lk}$, and
$(k,l)\mapsto(l,k)$ is a bijection of the index set.

Therefore
\[
  \bb_s(x,y)=\sum_e\sum_{k,l}
  \lambda\bigl(\sigma_s((x_e)_{kl})(y_e)_{kl}\bigr)
  =\sum_{e,k}\sum_{a=0}^{t-1}
  \sigma_s\bigl((x_e)_{k,a}\bigr)\,(y_e)_{k,t-1-a},
\]
because $\lambda(u^{a+b})=1$ iff $a+b=t-1$ and $0$ otherwise, where
$(x_e)_{k,a}$ denotes the $u^a$-coefficient of the $k$-th entry of $x_e$.

Reading the right-hand side in the flat coordinates
$u=\varphi_n(x)$, $v=\varphi_n(y)$, the summand pairs the coordinate
$(e,k,a)$ of $u$ (with $\sigma_s$ applied) with the coordinate
$(e,k,t-1-a)$ of $v$. In the ordering of Definition~\ref{def:gray} the
$a$-index runs inside a fixed $(e,k)$ block, so the Gram matrix is the identity
on the $(e,k)$-index tensored with $H_t$ on the $a$-index, i.e.
$J=I_{4n}\otimes H_t$.

The stated properties of $J$: $H_t$ is the reversal (Hankel) matrix, symmetric
with $H_t^2=I_t$; hence $J$ is symmetric with $J^2=I_{4n}\otimes H_t^2=I$,
non-singular. $J=I$ iff $H_t=I_t$, i.e. $t=1$. Finally $H_t$ is alternating
(zero diagonal, since $2a=t-1$ has no solution when $t$ is even) exactly when
$t$ is even, and $J$ inherits this.
\end{proof}

\begin{corollary}[the transported form is not Euclidean]\label{cor:noteuclid}
The transported form is Euclidean ($\langle u,v\rangle^{\widehat{}}_s=u\cdot v$)
if and only if $t=1$ and $r\mid s$; it is the plain
$\widehat\sigma_s$-Hermitian form if and only if $t=1$.
\end{corollary}

\begin{proof}
The Gram matrix of the Euclidean form is $I$; by
Theorem~\ref{thm:gram} the Gram matrix is $I_{4n}\otimes H_t$ with the
$\sigma_s$ twist, which equals $I$ iff $H_t=I_t$ (i.e. $t=1$) and
$\sigma_s=\id$ (i.e. $r\mid s$). Dropping the $\sigma_s$ condition gives the
plain Hermitian statement.
\end{proof}

\subsection{Duality and hull preservation}

\begin{theorem}[duality preservation]\label{thm:grayduality}
For every code $C\le\Rn$ and every $s$,
\[
  \varphi_n(C^{\perp_s})=\varphi_n(C)^{\hhatpers},
\]
where $D^{\hhatpers}=\{v:\langle u,v\rangle^{\widehat{}}_s=0\ \forall u\in D\}$.
\end{theorem}

\begin{proof}
$\varphi_n$ is bijective and $\langle\cdot,\cdot\rangle^{\widehat{}}_s$ is
defined as its transport, so for $y\in R^n$,
\[
  y\in C^{\perp_s}
  \iff \bb_s(x,y)=0\ \forall x\in C
  \iff \langle\varphi_nx,\varphi_ny\rangle^{\widehat{}}_s=0\ \forall x\in C
  \iff \varphi_n(y)\in\varphi_n(C)^{\hhatpers}. \qedhere
\]
\end{proof}

\begin{theorem}[hull preservation, with no principality hypothesis]
\label{thm:grayhull}
For every $C\le\Rn$ and every $s$,
\[
  \varphi_n\bigl(\Hull_s(C)\bigr)
  =\varphi_n(C)\cap\varphi_n(C)^{\hhatpers}.
\]
In particular the identity holds for the non-principal hulls of
Theorem~\ref{thm:nonprincipal}.
\end{theorem}

\begin{proof}
$\varphi_n$ is injective, hence preserves intersections:
$\varphi_n(C\cap C^{\perp_s})=\varphi_n(C)\cap\varphi_n(C^{\perp_s})$; now
apply Theorem~\ref{thm:grayduality}. No principality hypothesis is used at any
point: Definition~\ref{def:hull} is an intersection of $\F_q$-subspaces and
Theorems~\ref{thm:hullstruct} and~\ref{thm:grayduality} hold for arbitrary right
$\Rn$-submodules.
\end{proof}

\subsection{The image code}

\begin{proposition}[quasi-cyclicity]\label{prop:qc}
Let $D=\varphi_n(C)\subseteq\F_q^{\Nlen}$. Then $D$ is $\F_q$-linear of length
$\Nlen$ and dimension $\dim_{\F_q}C$, and it is \emph{quasi-cyclic of index
$4t$ and co-index $n$}: it is invariant under the rotation by $4t$
positions. It is not cyclic in general.
\end{proposition}

\begin{proof}
Linearity and dimension are Lemma~\ref{lem:philin}. Invariance under the
rotation by $4t$ is the shift identity $\varphi_nT=\widehat T\varphi_n$
together with $T$-invariance of the right ideal $C$. Non-cyclicity: the
rotation by one position does not in general preserve $D$; this is a
measurement on the enumerated families and we record it rather than claim it.
\end{proof}

\begin{remark}[what the map does to weights]\label{rem:weights}
Let $w_{\rm H}$ denote the Hamming weight on $R^n$ (counting nonzero
$R$-positions), $w_{\rm ent}$ the weight counting nonzero entries of the
$2\times2$ matrices, and $w_{\rm coef}$ the number of nonzero $\F_q$-coefficients.
Then
\[
  w_{\rm H}\le w_{\rm ent}\le w_{\rm coef}\le t\,w_{\rm ent},
  \qquad
  w_{\rm H}\le\wt_{\rm H}\varphi_n\le 4t\,w_{\rm H},
\]
and both bounds are sharp. In particular $\varphi_n$ \emph{is} an isometry from
the coefficient metric $d_{\rm coef}$ to the Hamming metric on
$\F_q^{\Nlen}$. The naive claim that $\varphi_n$ preserves the ring-Hamming
weight is false: a nonzero matrix with a single nonzero entry has
$w_{\rm H}=1$ but $\wt_{\rm H}\varphi_n\ge1$ with equality only in special
cases, and in general the two differ.
\end{remark}

\begin{remark}[no homogeneous-weight isometry]\label{rem:noisometry}
The homogeneous weight on $R=M_2(\F_q)$ takes the values $0$,
$1+\frac1{q^2-1}$ (rank one) and $1-\frac1{(q-1)^2(q+1)}$ (units). Since these
values are not integers, \emph{no} map into any Hamming space can be an isometry
for it. The chain of rings for which such isometries do exist
\cite{GreferathSchmidt1999GrayIsometry} consists of finite \emph{chain} rings,
and $R$ is not one. This places a hard limit on what any Gray map can achieve
here, and it is the reason Section~\ref{subsec:phys} recomputes distances in
physical coordinates rather than transporting them.
\end{remark}

\subsection{Every Galois hull is self-orthogonal}

This is the hinge of the paper. It holds in \emph{all three regimes} of
Definition~\ref{def:regimes}, not only the reflexive one.

\begin{theorem}[hull self-orthogonality]\label{thm:hullso}
For every $C\le\Rn$ and every $s$,
\[
  \varphi_n\bigl(\Hull_s(C)\bigr)
  \subseteq
  \varphi_n\bigl(\Hull_s(C)\bigr)^{\hhatpers}.
\]
No reflexivity hypothesis ($r\mid 2s$) and no principality hypothesis is
required.
\end{theorem}

\begin{proof}
Let $H=\Hull_s(C)=C\cap C^{\perp_s}$ and let $x,y\in H$. Then $x\in C$ and
$y\in C^{\perp_s}$. By the definition of $C^{\perp_s}$ we have
$\bb_s(c,y)=0$ for every $c\in C$; taking $c=x$ gives $\bb_s(x,y)=0$.
Transporting with $\varphi_n$ (Theorem~\ref{thm:gram}) gives
$\langle\varphi_nx,\varphi_ny\rangle^{\widehat{}}_s=0$.

Note what the argument uses and what it does not: it uses $x\in C$ and
$y\in C^{\perp_s}$, and it never interchanges the two slots. It therefore does
not need the orthogonality relation to be symmetric, which is exactly the
property that holds only in the reflexive regime
(Proposition~\ref{prop:regimes}). Nor does it use principality: $H$ is an
intersection of $\F_q$-subspaces, and Theorems~\ref{thm:hullstruct}
and~\ref{thm:grayduality} hold for arbitrary right $\Rn$-submodules.

Finally, note the direction of the statement: a nonzero hull dimension never
implies that $C$ itself is self-orthogonal.
\end{proof}

\subsection{Physical coordinates and the characteristic-two obstruction}
\label{subsec:phys}

For the Euclidean (CSS) route one wants the transported form to be the Euclidean
form. Corollary~\ref{cor:noteuclid} says this happens only for $t=1$ with
$r\mid s$; otherwise one looks for a linear change of coordinates.

\begin{definition}\label{def:congruence}
Assume $\sigma_s=\id$ (i.e. $r\mid s$), so that the transported form is the
bilinear form with matrix $J$. A \emph{congruence} is a matrix $\mathcal T$
with $\mathcal T\mathcal T^{\mathsf T}=J$; the associated coordinate map is
$\Psi(u)=u\,\mathcal T$.
\end{definition}

\begin{proposition}[the congruence and its cost]\label{prop:congruence}
Let $T$ be a $t\times t$ matrix with $T^{\mathsf T}H_tT=I_t$ and let
$\mathcal T=I_{4n}\otimes T$. Then $\mathcal T\mathcal T^{\mathsf T}=J$ (using
$J^2=I$ and $H_t^{\mathsf T}=H_t$), $\mathcal T$ is invertible, and $\Psi$
transports the transported form to the Euclidean form:
\[
  \Psi(u)\cdot\Psi(v)=u\,\mathcal T\mathcal T^{\mathsf T}v^{\mathsf T}
  =u\,J\,v^{\mathsf T}=\langle u,v\rangle^{\widehat{}}_s .
\]
Moreover $\mathcal T$ commutes with the block shift (so the quasi-cyclic
structure of index $4t$ is preserved), and $\mathcal T$ is monomial iff $t=1$.
\end{proposition}

\begin{proof}
Immediate from $J=I_{4n}\otimes H_t$, $H_t^{\mathsf T}=H_t$, $J^2=I$ and the
definitions; $\mathcal T$ is block-diagonal in the flat ordering, hence commutes
with the block shift. Monomiality: for $t=1$, $H_1=(1)$ and $T$ may be taken to
be $(1)$; for $t\ge2$, $T^{\mathsf T}H_tT=I_t$ forces $T$ to have at least two
nonzero entries in some row, so $\mathcal T$ is not monomial.
\end{proof}

\begin{theorem}[characteristic-two obstruction]\label{thm:obstruction}
Let $p=2$ and let $t$ be even. Then no matrix $T$ with
$T^{\mathsf T}H_tT=I_t$ exists. Consequently, for every such $(q,t)$ and
\emph{every} code, the Euclidean/CSS route is not available: no linear change of
coordinates converts the transported form into the Euclidean form.
\end{theorem}

\begin{proof}
If $t$ is even then $H_t$ is alternating: its diagonal entries are
$(H_t)_{aa}=1$ iff $2a=t-1$, which has no integer solution when $t$ is even;
and $H_t$ is symmetric. A symmetric matrix congruent to the identity over a
field of characteristic $2$ cannot be alternating, because an alternating
matrix $M$ satisfies $x^{\mathsf T}Mx=0$ for all $x$ whereas
$x^{\mathsf T}Ix=|x|^2$ is not identically zero. Hence no such $T$ exists.
\end{proof}

\begin{remark}[the hypotheses are exact]\label{rem:obstructionexact}
Theorem~\ref{thm:obstruction} is stated for $p=2$ \emph{and} $t$ even, and it
must not be extended. Congruences do exist for $t=1$, $t=3$ and $t=5$ at
$q=4,8$, and one exists at $(q,t)=(9,2)$. What the obstruction forbids is
exactly the pair ``characteristic two, $t$ even''; it is a statement about the
form $J$, not about any code. Where a congruence exists but is non-monomial
(Proposition~\ref{prop:congruence}), it must not be treated as an isometry:
measured instances show that the minimum distance of a code \emph{does} change
under $\Psi$ (3 of 5 codes at $(9,2,2)$), so every distance must be recomputed
in physical coordinates.
\end{remark}

\begin{theorem}[number of physical coordinates]\label{thm:nphys}
A quantum code produced from a code $C\le\Rn$ by any of the routes of
Section~\ref{sec:quantum} has length $\Nlen=4tn$ over the alphabet $\F_q$.
\end{theorem}

\begin{proof}
One qudit of dimension $q$ carries one $\F_q$-coordinate, so the number of
physical qudits equals the $\F_q$-length of the classical code whose symplectic
group lives over $\F_q^{2\Nlen}$: by Lemma~\ref{lem:basic}(c),
$\dim_{\F_q}R^n=\Nlen=4tn$. There is no doubling to $2\Nlen$: $\F_q^{2\Nlen}$
is the symplectic \emph{bookkeeping} space in which the stabilizer group sits,
not a set of qudits. Nor is there a substitution of $\F_q$ by $\F_p$: the
alphabet of the quantum code is $\F_q$ throughout.
\end{proof}

\subsection{Worked examples IV: the Gray map}

\begin{example}[the coefficient expansion and the Gram matrix, at $t=2$]
\label{ex:gray}
Let $q=4$, $t=2$, $n=1$, so $\Nlen=4\cdot2\cdot1=8$. Take
\[
  X=\begin{pmatrix}1+\alpha u&\beta u\\ 0&1\end{pmatrix}\in R=M_2(A_2),
  \qquad \alpha,\beta\in\F_4 .
\]
Then $\varphi_A(1+\alpha u)=(1,\alpha)$, $\varphi_A(\beta u)=(0,\beta)$,
$\varphi_A(0)=(0,0)$, $\varphi_A(1)=(1,0)$, so
\[
  \varphi_R(X)=\bigl(1,\alpha,\;0,\beta,\;0,0,\;1,0\bigr)\in\F_4^8 ,
\]
in the flat ordering of Definition~\ref{def:gray}: the four $2\times2$ entries
in reading order, each expanded into its $t=2$ coefficients.

For $t=2$ the Hankel matrix is
\[
  H_2=\begin{pmatrix}0&1\\1&0\end{pmatrix},
  \qquad\text{so}\qquad
  J=I_{4n}\otimes H_2=I_4\otimes H_2 ,
\]
an $8\times8$ matrix which is symmetric, satisfies $J^2=I_8$, is \emph{not}
the identity (in accordance with Corollary~\ref{cor:noteuclid}, since
$t\ne1$), and \emph{is} alternating, since $(H_2)_{aa}=1$ would require
$2a=1$ and $p=2$. The transported form is therefore
\[
  \langle u,v\rangle^{\widehat{}}_s=u^{[s]}Jv^{\mathsf T}
  =\sum_{i=0}^{3}\bigl(\sigma_s(u_{2i})v_{2i+1}
    +\sigma_s(u_{2i+1})v_{2i}\bigr),
\]
which pairs the two $u$-coefficients inside each of the four matrix entries
and nothing else. Note that it is not the dot product $u\cdot v$: the
coordinate $u_0$ is paired with $v_1$, not with $v_0$.
\end{example}

\begin{example}[the characteristic-two obstruction, at $(q,t)=(4,2)$]
\label{ex:obstruction}
Let $q=4$ (so $p=2$) and $t=2$. Then $H_2=\begin{pmatrix}0&1\\1&0\end{pmatrix}$
has zero diagonal, and for every $x=(x_0,x_1)^{\mathsf T}\in\F_4^2$,
\[
  x^{\mathsf T}H_2x=x_0x_1+x_1x_0=2x_0x_1=0
\]
in characteristic $2$: $H_2$ is alternating. For any $2\times2$ matrix $T$ over
$\F_4$ and any $x$,
\[
  x^{\mathsf T}(T^{\mathsf T}H_2T)x=(Tx)^{\mathsf T}H_2(Tx)=0,
\]
so $T^{\mathsf T}H_2T$ is alternating for \emph{every} $T$. But $I_2$ is not
alternating, since $e_1^{\mathsf T}I_2e_1=1\ne0$. Hence no $T$ with
$T^{\mathsf T}H_2T=I_2$ exists; the search over all $4^4=256$ matrices
$T\in M_2(\F_4)$ confirms this exhaustively, and Theorem~\ref{thm:obstruction}
proves it for every even $t$ in characteristic $2$.

The hypotheses are exact, and the neighbouring cases are different: for
$t=1,3,5$ at $q=4,8$ a congruence is \emph{constructed} (symmetric elimination
with Frobenius scaling, backtracking over pivots) and verified by direct
multiplication, and one exists at $(q,t)=(9,2)$ as well --- there $p=3$ and the
argument above breaks because $2x_0x_1$ need not vanish. Where a congruence
exists but is non-monomial (Proposition~\ref{prop:congruence}), it must not
be treated as an isometry: at $(9,2,2)$ the minimum distance changed under
$\Psi$ for $3$ of the $5$ codes tested, so every distance in this paper is
recomputed in physical coordinates.
\end{example}

% =====================================================================
\section{Quantum error-correcting code constructions}
\label{sec:quantum}
% =====================================================================

Nothing in this section is imported: the stabilizer formalism is derived from
the commutation relation for the $q$-ary Weyl operators, so that every
hypothesis is visible where it is used.

\subsection{The $q$-ary Weyl operators}

Let $\Nlen\ge1$ and let $\mathcal H=(\mathbb C^q)^{\otimes\Nlen}$ with basis
$\{|v\rangle:v\in\F_q^{\Nlen}\}$. For $a,b\in\F_q^{\Nlen}$ define
\begin{equation}\label{eq:weyl}
  X(a)|v\rangle=|v+a\rangle,\qquad
  Z(b)|v\rangle=\chi\bigl(\Tr_{\F_q/\F_p}(b\cdot v)\bigr)|v\rangle,
  \qquad \chi(x)=e^{2\pi i x/p},
\end{equation}
and $E(a|b)=X(a)Z(b)$. Write
\begin{equation}\label{eq:omega}
  \omega\bigl((a|b),(a'|b')\bigr)=a\cdot b'-a'\cdot b\ \in\F_q,
  \qquad
  \swt(a|b)=\#\{i:(a_i,b_i)\ne(0,0)\}.
\end{equation}

\begin{lemma}[commutation]\label{lem:Q1}
$X(a)Z(b)=\chi(\Tr(a\cdot b))\,Z(b)X(a)$. Consequently two Weyl operators
$E(a|b)$ and $E(a'|b')$ commute if and only if
$\Tr_{\F_q/\F_p}\bigl(\omega((a|b),(a'|b'))\bigr)=0$.
\end{lemma}

\begin{proof}
$Z(b)X(a)|v\rangle=Z(b)|v+a\rangle
=\chi(\Tr(b\cdot(v+a)))|v+a\rangle
=\chi(\Tr(b\cdot a))\,X(a)Z(b)|v\rangle$, which gives the first claim after
swapping the roles. For the second, expand
$E(a|b)E(a'|b')=\chi(\Tr(a'\cdot b))E(a+a'|b+b')$ and
$E(a'|b')E(a|b)=\chi(\Tr(a\cdot b'))E(a+a'|b+b')$ and compare the phases.
\end{proof}

Lemma~\ref{lem:Q1} is the reason commutation is a condition over $\F_p$ (a
trace), even though the code is defined over $\F_q$. This has a consequence that
is easy to get wrong; see Theorem~\ref{thm:fpbasis}.

\begin{lemma}[stabilizer lemma]\label{lem:Q2}
Let $\mathcal C\subseteq\F_q^{2\Nlen}$ be an additive (that is, $\F_p$-linear)
subgroup and let $S=\{E(v):v\in\mathcal C\}$ be the group it generates. Then:
\begin{enumerate}[label=\textnormal{(\alph*)},leftmargin=*]
\item $S$ is abelian iff $\Tr(\omega(x,y))=0$ for all $x,y\in\mathcal C$;
\item the common eigenspace of $S$ has dimension $q^{\Nlen}/|\mathcal C|$;
  writing $K=\log_q(q^{\Nlen}/|\mathcal C|)$ one has
  $K=\Nlen-\log_q|\mathcal C|$;
\item with $\mathcal C_0=\mathcal C\cap\mathcal C^{\perp_{\Tr\circ\omega}}$,
  the code detects every error of symplectic weight $<d$ and corrects every
  error of weight $\le\lfloor(d-1)/2\rfloor$, where
  \[
    d=\min\{\swt(v):v\in\mathcal C^{\perp_{\Tr\circ\omega}}\setminus
    \mathcal C_0\};
  \]
\item (Knill--Laflamme, group form) the code corrects an error set
  $\mathcal E$ iff for all $E,F\in\mathcal E$ either $E^\dagger F\in\pm S$ or
  some element of $S$ anti-commutes with $E^\dagger F$.
\end{enumerate}
\end{lemma}

\begin{proof}
(a) is Lemma~\ref{lem:Q1}. (b) The operators $\{E(v)\}$ for $v$ ranging over a
set of representatives of $\F_q^{2\Nlen}/\mathcal C$ are pairwise distinct and
there are $q^{2\Nlen}/|\mathcal C|$ of them; the space decomposes into
$|\mathcal C|$ joint eigenspaces permuted transitively, whence the dimension.
(c) An error of symplectic weight $w$ is a linear combination of Weyl operators
of symplectic weight $\le w$; such an operator is undetected iff it lies in the
centralizer of $S$ and outside $S$ up to phases, i.e. iff its label lies in
$\mathcal C^{\perp_{\Tr\circ\omega}}\setminus\mathcal C_0$. (d) is the standard
translation of the Knill--Laflamme conditions into the group language.
\end{proof}

We write $[[\Nlen,K,d]]_q$ for a $q$-ary quantum code of length $\Nlen$,
dimension $q^K$ and distance $d$, and $[[\Nlen,K,d;c]]_q$ for an
\emph{entanglement-assisted} code consuming $c$ maximally entangled $q$-ary
pairs.

\begin{theorem}[commutation must be tested over $\F_p$]\label{thm:fpbasis}
Let $\mathcal C$ be $\F_q$-linear. Then
$\mathcal C^{\perp_{\Tr\circ\omega}}=\mathcal C^{\perp_\omega}$, where
$\perp_\omega$ is the $\F_q$-symplectic orthogonal. However, testing
commutation only on an $\F_q$-basis of $\mathcal C$ is \textbf{invalid}:
there exists a code --- the $[[16,4,3]]_4$ code of
Section~\ref{sec:cert} --- and a weight-one operator $(0|e_0)$ which commutes
with all $12$ elements of an $\F_q$-generating set of the stabilizer yet not
with all $24$ elements of an $\F_p$-basis, because
$\Tr_{\F_4/\F_2}(1)=0$. The naive test reports $d=1$ instead of $d=3$.
\end{theorem}

\begin{proof}
The equality of the two orthogonals is the non-degeneracy of the trace form: if
$\Tr(\omega(x,c))=0$ for all $c\in\mathcal C$ then, by $\F_q$-linearity of
$\mathcal C$ and non-degeneracy of $\Tr_{\F_q/\F_p}$, $\omega(x,c)=0$ as an
element of $\F_q$.

For the counterexample, let $\mathcal C$ be the stabilizer of the
$[[16,4,3]]_4$ code and let $v=(0|e_0)$. Then
$\omega(v,(a|b))=-a_0$, so $v$ commutes with $(a|b)$ iff
$\Tr_{\F_4/\F_2}(a_0)=0$. The set of $a_0\in\F_4$ with zero absolute trace is
$\{0,1\}$, a $1$-dimensional $\F_2$-subspace of $\F_4$ that is \emph{not} an
$\F_4$-subspace. Hence commutation holds on an $\F_4$-generating set only if
all their first-slot symbols happen to lie in $\{0,1\}$ --- which is the
observed degeneracy --- but fails on a full $\F_2$-basis. The correct criterion
is $\omega(v,c)=0$ as an element of $\F_q$ for all $c\in\mathcal C$,
equivalently commutation with every element of an $\F_p$-basis.
\end{proof}

\subsection{Three routes}

Throughout this subsection $D\subseteq\F_q^{\Nlen}$ is $\F_q$-linear of
dimension $k$, $D^{\perp_E}$ is its Euclidean dual and
$H_E=\Hull_E(D)=D\cap D^{\perp_E}$ with $h_E=\dim H_E$. We write
\begin{equation}\label{eq:V}
  V=D\oplus D=\{(a|b):a,b\in D\}\subseteq\F_q^{2\Nlen},
  \qquad
  V^{\omegaperp}=D^{\perp_E}\oplus D^{\perp_E},
\end{equation}
so that $\rad V=V\cap V^{\omegaperp}=H_E\oplus H_E$.

\begin{theorem}[route R1: the Galois-native twisted embedding]
\label{thm:R1}
Let $D\subseteq\F_q^{\Nlen}$ and put $T_s(u)=(u^{[s]}J)^{\mathsf T}$ and
\[
  \mathcal C_T=\{(u|T_su):u\in D\}\subseteq\F_q^{2\Nlen}.
\]
Then for all $u,u'\in D$,
\[
  \omega\bigl((u|T_su),(u'|T_su')\bigr)
  =\langle u',u\rangle^{\widehat{}}_s-\langle u,u'\rangle^{\widehat{}}_s .
\]
Consequently, if $D$ is self-orthogonal for the transported form then
$\mathcal C_T$ is isotropic, the operators $E(u|T_su)$ pairwise commute, and
each has order $p$, so that Lemma~\ref{lem:Q2} gives a quantum code
\[
  [[\Nlen,\Nlen-k,d]]_q,\qquad
  d=\min\{\swt(v):v\in\mathcal C_T^{\perp_{\Tr\circ\omega}}\setminus
  \mathcal C_T\}.
\]
No coordinate transformation and no reflexivity hypothesis is used.
\end{theorem}

\begin{proof}
$(T_su)_i=\sum_j\sigma_s(u_j)J_{ji}$, hence
$u\cdot(T_su')=\sum_{i,j}u_i\sigma_s(u'_j)J_{ji}
=\langle u',u\rangle^{\widehat{}}_s$ after relabelling $i\leftrightarrow j$ and
using that $J$ is symmetric; similarly
$u'\cdot(T_su)=\langle u,u'\rangle^{\widehat{}}_s$. Subtracting gives the
identity. If $D$ is self-orthogonal both terms vanish, so $\omega=0$ on
$\mathcal C_T$ and Lemma~\ref{lem:Q2}(a) applies. Each operator has order $p$
because $u\cdot(T_su)=\langle u,u\rangle^{\widehat{}}_s=0$. Finally
$|\mathcal C_T|=|D|=q^k$, so $K=\Nlen-k$ by Lemma~\ref{lem:Q2}(b).
\end{proof}

\begin{theorem}[the R1 distance cap]\label{thm:R1cap}
For every $b\in D^{\perp_E}$ the vector $(0|b)$ lies in
$\mathcal C_T^{\perp_\omega}$. Consequently
\[
  d_{R1}\le\min\wt\bigl(D^{\perp_E}\bigr).
\]
\end{theorem}

\begin{proof}
$\omega((0|b),(u|T_su))=0\cdot(T_su)-u\cdot b=-u\cdot b$, which vanishes for
all $u\in D$ exactly when $b\in D^{\perp_E}$. Hence
$(0|b)\in\mathcal C_T^{\perp_\omega}$ and $\swt(0|b)=\wt(b)$.
\end{proof}

Theorem~\ref{thm:R1cap} is why R1 is never the route of choice for distance. It
remains the only route attached to the Galois structure when the obstruction of
Theorem~\ref{thm:obstruction} applies.

\begin{theorem}[route R2: Euclidean/CSS]\label{thm:R2}
Let $D'\subseteq\F_q^{\Nlen}$ be $\F_q$-linear of dimension $k$.
\begin{enumerate}[label=\textnormal{(\alph*)},leftmargin=*]
\item If $D'\subseteq D'^{\perp_E}$ (Euclidean self-orthogonal) then
  $\mathcal C=D'\oplus D'$ is isotropic and yields
  $[[\Nlen,\Nlen-2k,d]]_q$ with $d=\min\wt(D'^{\perp_E}\setminus D')$.
\item If $D'^{\perp_E}\subseteq D'$ (dual-containing) then
  $\mathcal C=D'^{\perp_E}\oplus D'^{\perp_E}$ is isotropic and yields
  $[[\Nlen,2k-\Nlen,d]]_q$ with $d=\min\wt(D'\setminus D'^{\perp_E})$.
\end{enumerate}
\end{theorem}

\begin{proof}
Take $\mathcal C=A\oplus B$ with $A,B\subseteq\F_q^{\Nlen}$. Then
$\omega((a|0),(0|b))=a\cdot b$, so $\mathcal C$ is isotropic iff
$A\subseteq B^{\perp_E}$; moreover $\mathcal C^{\omegaperp}
=B^{\perp_E}\oplus A^{\perp_E}$ and $\dim\mathcal C=\dim A+\dim B$, giving
$K=\Nlen-\dim A-\dim B$ by Lemma~\ref{lem:Q2}(b). Taking $A=B=D'$ gives (a) and
$A=B=D'^{\perp_E}$ gives (b).

For the distance: a vector $(a|b)\in\mathcal C^{\omegaperp}\setminus\mathcal C$
has $a\notin A$ or $b\notin B$, and
$\swt(a|b)\ge\max(\wt a,\wt b)$; conversely $(a|0)$ realises $\wt(a)$ for
$a\in A^{\perp_E}\setminus A$ and $(0|b)$ realises $\wt(b)$ for
$b\in B^{\perp_E}\setminus B$. Hence the symplectic minimum equals the stated
classical minimum.
\end{proof}

\begin{theorem}[route R3: entanglement-assisted]\label{thm:R3}
Let $D\subseteq\F_q^{\Nlen}$ be \emph{any} $\F_q$-linear code of dimension $k$,
and put $\mathcal C=V=D\oplus D$ as in \eqref{eq:V}. Then
\[
  [[\Nlen,\Nlen-k-h_E,\ d\ ;\ k-h_E]]_q,\qquad
  d=\min\wt\bigl(D^{\perp_E}\setminus\Hull_E(D)\bigr).
\]
For $h_E=k$ (i.e. $D$ Euclidean self-orthogonal) this degenerates to $c=0$ and
reproduces Theorem~\ref{thm:R2}(a).
\end{theorem}

\begin{proof}
The isotropic part of $V$ is
$\rad V=V\cap V^{\omegaperp}=H_E\oplus H_E$, of dimension $2h_E$. The
commutativity defect is therefore
$2c=\dim V-\dim\rad V=2k-2h_E$, i.e. $c=k-h_E$ maximally entangled $q$-ary
pairs.

With entanglement available, the \emph{whole} group $V$ becomes measurable.
Decompose $\F_q^{2\Nlen}$ as $\rad V\oplus\bigoplus_i(\F_qe_i\oplus\F_qf_i)$
with $\{e_i,f_i\}$ hyperbolic pairs and $V=\rad V\oplus\bigoplus_i\F_qe_i$.
For $v=s+\sum_i(\alpha_ie_i+\beta_if_i)$ with $s\in\rad V$, the element $s$
stabilizes the codespace outright, and each hyperbolic pair is measured with the
help of one ebit: on the maximally entangled state $|\Phi\rangle$ one has
$(A\otimes I)|\Phi\rangle=(I\otimes A^{\mathsf T})|\Phi\rangle$, so
$E(e_i)\otimes I$ and $I\otimes E(f_i)^{\mathsf T}$ act identically and the
pair is simultaneously measurable. Hence $v\otimes I$ fixes the codespace.

Consequently an error $F$ is undetected exactly when it commutes with all of $V$
\textbf{and} lies outside $V$ --- note: outside $V$, not merely outside
$\rad V$ --- i.e. when its label lies in $V^{\omegaperp}\setminus V$. By
\eqref{eq:V} and the argument of Theorem~\ref{thm:R2}, the symplectic minimum
over $V^{\omegaperp}\setminus V$ equals
$\min\wt(D^{\perp_E}\setminus\Hull_E(D))$. The dimension is
$K=\Nlen-\dim V+c=\Nlen-2k+(k-h_E)=\Nlen-k-h_E$.

For $h_E=k$ we get $c=0$, $K=\Nlen-2k$ and
$D^{\perp_E}\setminus\Hull_E(D)=D^{\perp_E}\setminus D$, which is
Theorem~\ref{thm:R2}(a).
\end{proof}

\begin{remark}[field discipline]
$\mathcal C=D\oplus D$ is $\F_q$-linear, so by
Theorem~\ref{thm:fpbasis} the trace-symplectic and the $\F_q$-symplectic
orthogonals coincide and the entanglement count $c=k-h_E$ is an $\F_q$-quantity.
For the \emph{semilinear} group $\mathcal C_T$ of R1 the standard $q$-ary
entanglement formula does not apply unchanged, because $\mathcal C_T$ is only
$\F_p$-linear; we do not use it.
\end{remark}

\subsection{The parameter theorem}

\begin{theorem}[parameters]\label{thm:parameters}
Let $C\le\Rn$, $D=\varphi_n(C)$, $\Nlen=4tn$, $k=\dim_{\F_q}D$ and
$h_E=\dim(D\cap D^{\perp_E})$. Then:
\begin{enumerate}[label=\textnormal{(\alph*)},leftmargin=*]
\item \textbf{(R1)} if $D\subseteq D^{\hhatpers}$ then $[[\Nlen,\Nlen-k,d]]_q$
  with $d\le\min\wt(D^{\perp_E})$;
\item \textbf{(R2)} if $D'$ is Euclidean self-orthogonal of dimension $k$ then
  $[[\Nlen,\Nlen-2k,d]]_q$, and if $D'$ is dual-containing then
  $[[\Nlen,2k-\Nlen,d]]_q$; here $D'=D$ when $t=1$ and $r\mid s$, and
  $D'=\Psi(D)$ when a congruence as in Proposition~\ref{prop:congruence}
  exists;
\item \textbf{(R3)} for arbitrary $D$,
  $[[\Nlen,\Nlen-k-h_E,d;k-h_E]]_q$.
\end{enumerate}
In all cases, writing $G$ for the group ($\mathcal C_T$ for R1, $V$ for R2/R3),
\[
  d=\min\{\swt(v):v\in G^{\omegaperp}\setminus G\}.
\]
All dimensions are over $\F_q$, the length is $\Nlen=4tn$, and the alphabet is
$\F_q$.
\end{theorem}

\begin{proof}
Combine Theorem~\ref{thm:nphys} for the length, Theorems~\ref{thm:R1},
\ref{thm:R2} and~\ref{thm:R3} for the three routes, and
Lemma~\ref{lem:Q2}(c) for the distance formula. Note that the reduction of the
symplectic minimum to a classical minimum is proved for R2 and R3 only; for R1
only the upper bound of Theorem~\ref{thm:R1cap} is available, and we state it as
such.
\end{proof}

\begin{remark}[which route to use]
R2 is the primary route whenever the transported form is already Euclidean or a
congruence exists: it is the only route whose distance is not capped by
$\min\wt(D^{\perp_E})$, and it consumes no entanglement. R1 is the
Galois-native route for $t\ge2$ when no congruence exists --- precisely the case
of the non-principal hull witnesses --- accepting the cap of
Theorem~\ref{thm:R1cap}. R3 converts \emph{any} image code, including an LCD one
($h_E=0$), into a quantum code, at the price of entanglement.
\end{remark}

\subsection{Worked examples V: the $\F_p$-basis subtlety}

\begin{example}[why commutation must be tested over $\F_p$]\label{ex:fpbasis}
Let $q=4$, $\Nlen=16$, and consider the $[[16,4,3]]_4$ code produced at
$(q,t,n)=(4,1,4)$ by the self-orthogonal route: the image code
$D\subseteq\F_4^{16}$ has $k=6$, $h_E=6$, so $K=16-2\cdot6=4$. The stabilizer
group is $\mathcal C=D\oplus D$, of $\F_4$-dimension $12$.

Take the weight-one operator $E(0|e_0)$ labelled by $v=(0|e_0)$, where
$e_0=(1,0,\dots,0)$. For $(a|b)\in\mathcal C$,
\[
  \omega\bigl(v,(a|b)\bigr)=0\cdot b-a\cdot e_0=-a_0 .
\]
So $E(v)$ commutes with $E(a|b)$ iff $\Tr_{\F_4/\F_2}(a_0)=0$. Now
\[
  \Tr_{\F_4/\F_2}(1)=1+1^2=0,\qquad \Tr_{\F_4/\F_2}(\omega)=\omega+\omega^2=1,
\]
so the set of $a_0\in\F_4$ with vanishing trace is $\{0,1\}$: a
$1$-dimensional $\F_2$-subspace of $\F_4$ that is \emph{not} an
$\F_4$-subspace. Consequently:
\begin{itemize}[leftmargin=*]
\item testing commutation on the $12$ elements of an $\F_4$-generating set of
  $\mathcal C$ reports that $v$ commutes with everything, and the naive
  distance comes out as $d=1$;
\item testing it on the $24$ elements of an $\F_2$-basis detects that some
  $(a|b)$ has $a_0=\omega$, so $v$ does \emph{not} commute, and the correct
  distance $d=3$ is obtained.
\end{itemize}
By Theorem~\ref{thm:fpbasis} the two orthogonals
$\mathcal C^{\perp_{\Tr\circ\omega}}$ and $\mathcal C^{\perp_\omega}$
\emph{coincide} for an $\F_q$-linear $\mathcal C$; what fails is not the
orthogonal but the \emph{test}. Every distance reported in this paper is
computed with the $\F_p$-basis test.
\end{example}

\begin{example}[the three routes on one family]\label{ex:routes}
Take $(q,t,n)=(9,1,6)$, so $\Nlen=24$. The search produces, among others, two
image codes $D\subseteq\F_9^{24}$ with very different hulls.

\emph{First,} $k=\dim D=8$ and $h_E=8$, i.e. $D\subseteq D^{\perp_E}$: the
code is Euclidean self-orthogonal. Route R2(a)
(Theorem~\ref{thm:R2}) applies with $A=B=D$ and gives
\[
  K=\Nlen-2k=24-16=8,\qquad
  d=\min\wt\bigl(D^{\perp_E}\setminus D\bigr)=4,\qquad
  [[24,8,4]]_9 .
\]
Feeding the same $D$ to route R3 instead gives $c=k-h_E=0$ and
$K=\Nlen-k-h_E=8$: the same code, which is the degeneracy stated at the end of
Theorem~\ref{thm:R3}. Note $D^{\perp_E}\setminus D=D^{\perp_E}\setminus
\Hull_E(D)$ here, since $\Hull_E(D)=D$; so the coset-distance theorem
(Theorem~\ref{thm:cosetdistance}) is the same statement for R2 and R3 in this
case.

\emph{Second,} $k=20$ and $h_E=2$, so $D$ is neither self-orthogonal nor
dual-containing and R2 does not apply at all. Route R3 gives
\[
  c=k-h_E=18,\qquad K=\Nlen-k-h_E=24-20-2=2,\qquad
  [[24,2,10;18]]_9 .
\]
This is the largest-distance code found anywhere in the search, and it costs
$18$ ebits out of a length of $24$.

The two rows are the two ends of the trade-off made explicit by
Theorem~\ref{thm:feasibility}: with $d\ge4$ and $K\ge2$ one needs
$2k\le\Nlen+c_0-K_0$, so a large $k$ (which is what makes $D^{\perp_E}$ small
and hence $d$ large) forces a large entanglement budget $c_0$. The
stabilizer code pays nothing and reaches $d=4$; the entanglement-assisted code
reaches $d=10$ and pays $c=18$.
\end{example}

% =====================================================================
\section{Distance certification and quantum parameter analysis}
\label{sec:cert}
% =====================================================================

\subsection{The coset-distance theorem}

The distance formulas of Theorems~\ref{thm:R2} and~\ref{thm:R3} look like
minima over a set difference. The following theorem says that they are
something better: the minimum weight of an ordinary linear code.

\begin{theorem}[coset-distance characterisation]\label{thm:cosetdistance}
Let $D\subseteq\F_q^{\Nlen}$ be $\F_q$-linear, $H_E=D\cap D^{\perp_E}$ and
$K=\Nlen-k-h_E$. Then, for routes R2 and R3,
\[
  d=\min\bigl\{\wt(x):x\in D^{\perp_E}\setminus\Hull_E(D)\bigr\}
\]
is the \emph{minimum weight of the quotient code}
\[
  D^{\perp_E}/\Hull_E(D),
\]
a linear code of length $\Nlen$ and dimension
$\dim D^{\perp_E}-\dim\Hull_E(D)=(\Nlen-k)-h_E=K$. Equivalently, $d$ is the
minimum weight of the set of nonzero cosets of $\Hull_E(D)$ in $D^{\perp_E}$.
In particular, for any integer $\delta\ge1$,
\begin{equation}\label{eq:criterion}
  d\ge\delta
  \iff
  \text{every }x\in D^{\perp_E}\text{ with }\wt(x)<\delta
  \text{ lies in }\Hull_E(D).
\end{equation}
\end{theorem}

\begin{proof}
By Theorem~\ref{thm:parameters}, for R2 and R3
$d$ is the minimum of $\swt$ over $G^{\omegaperp}\setminus G$ with $G=V$.
Since $\swt(a|b)\ge\max(\wt a,\wt b)$ and since
$(a|0)\in V^{\omegaperp}\setminus V$ whenever $a\in D^{\perp_E}\setminus D$,
that symplectic minimum equals
$\min\{\wt(x):x\in D^{\perp_E}\setminus D\}$; and for $x\in D^{\perp_E}$ one
has $x\in D$ iff $x\in D\cap D^{\perp_E}=H_E$, so
$D^{\perp_E}\setminus D=D^{\perp_E}\setminus H_E$ as sets.

Now, the quotient group $D^{\perp_E}/H_E$ is an $\F_q$-vector space of
dimension $\dim D^{\perp_E}-\dim H_E=(\Nlen-k)-h_E=K$, and its nonzero elements
are exactly the cosets $x+H_E$ with $x\in D^{\perp_E}\setminus H_E$. The weight
of a coset is well defined as
$\wt(x+H_E)=\min\{\wt(y):y\in x+H_E\}$, and its minimum over nonzero cosets is
exactly $\min\{\wt(x):x\in D^{\perp_E}\setminus H_E\}$: indeed
$\wt(x+H_E)\le\wt(x)$ gives one inequality, and conversely any $y\in x+H_E$
with $y\ne0$ may be substituted for $x$. This is the assertion. The criterion
\eqref{eq:criterion} is the contrapositive restatement.
\end{proof}

\begin{remark}
For route R1 only the upper bound of Theorem~\ref{thm:R1cap} is available, and
Theorem~\ref{thm:cosetdistance} is \emph{not} claimed for R1: the group
$\mathcal C_T$ is semilinear, the set difference is not a quotient by an
$\F_q$-subspace of the ambient space in the same way, and we do not use it.
\end{remark}

\subsection{The coset Singleton bound}

\begin{theorem}[coset Singleton bound]\label{thm:cosetsingleton}
Let $D\subseteq\F_q^{\Nlen}$ be $\F_q$-linear and let $K=\Nlen-k-h_E$ be the
dimension of the quantum code produced by route R2 or R3. Then
\begin{equation}\label{eq:cosetsingleton}
  K+d\le \Nlen+1 .
\end{equation}
\end{theorem}

\begin{proof}
Consider the quotient code $\mathcal Q=D^{\perp_E}/\Hull_E(D)$ of
Theorem~\ref{thm:cosetdistance}: it has length $\Nlen$, dimension $K$ and
minimum weight $d$. Let $S\subseteq\{1,\dots,\Nlen\}$ be any set of
$\Nlen-d+1$ coordinates and let $\pi_S$ be the projection onto those
coordinates. If $\pi_S$ were not injective on $\mathcal Q$, there would be a
nonzero coset $x+H_E$ with $\pi_S(x+H_E)=0$, i.e. an element
$y\in x+H_E$ supported outside $S$, i.e. supported on the remaining
$d-1$ coordinates. Then $\wt(y)\le d-1<d$, contradicting the definition of $d$
as the minimum weight of a nonzero element of $\mathcal Q$. Hence $\pi_S$ is
injective, so $q^K=|\mathcal Q|\le q^{\Nlen-d+1}$ and therefore
$K\le\Nlen-d+1$.
\end{proof}

\begin{proposition}[the coset Singleton bound is new]\label{prop:newsbound}
The bound \eqref{eq:cosetsingleton} is not the entanglement-assisted Singleton
bound, and neither implies the other.
\end{proposition}

\begin{proof}
The entanglement-assisted Singleton bound reads
$\Nlen-K+c\ge 2(d-1)$, i.e. $K\le\Nlen+c-2d+2$. For the code
$[[20,2,10;18]]_4$ constructed in Section~\ref{sec:cmp} one has $\Nlen=20$,
$K=2$, $d=10$, $c=18$: the coset Singleton bound gives $2\le 20-10+1=11$, which
holds, while the EA Singleton bound gives
$2\le 20+18-20+2=20$, which is much weaker. Conversely, for a code with small
$c$ and small $d$, say $[[24,20,2;0]]_4$ ($\Nlen=24$, $K=20$, $d=2$, $c=0$), the
coset Singleton bound gives $20\le 24-2+1=23$, while the EA Singleton bound
gives $20\le 24+0-4+2=22$ --- the EA bound is now the stronger one. Hence
neither implies the other.
\end{proof}

\subsection{The feasibility inequality}

\begin{theorem}[feasibility inequality]\label{thm:feasibility}
Suppose a code produced by route R3 from $D\subseteq\F_q^{\Nlen}$ has
$d\ge4$, $c\le c_0$ and $K\ge K_0$. Then
\begin{equation}\label{eq:feasibility}
  2k\le \Nlen+c_0-K_0 .
\end{equation}
\end{theorem}

\begin{proof}
By Theorem~\ref{thm:R3}, $c=k-h_E$ and $K=\Nlen-k-h_E$. The hypotheses
$c\le c_0$ and $K\ge K_0$ read $h_E\ge k-c_0$ and
$h_E\le \Nlen-k-K_0$. Hence $k-c_0\le \Nlen-k-K_0$, i.e.
$2k\le\Nlen+c_0-K_0$.
\end{proof}

\begin{remark}[what the inequality does and does not say]
The hypothesis $d\ge4$ enters only through the \emph{motivation} for imposing
the constraints, not through the algebra: the implication
``$c\le c_0$ and $K\ge K_0$ $\Rightarrow$ \eqref{eq:feasibility}'' holds for
every code, whatever its distance. Inequality \eqref{eq:feasibility} is one-way:
it is a necessary condition, not a sufficient one, and no converse is claimed.

Its content is the trade-off behind the entanglement budget: cheap entanglement
means a large Euclidean hull, and a large Euclidean hull means few logical
qudits. Combined with the criterion \eqref{eq:criterion}, it converts the
question ``is there a family with $d\ge4$, bounded $c$ and $K\ge2$?'' into two
explicit conditions on $D$: (i) $k$ must be close to $\Nlen/2$, and (ii) every
vector of weight at most $3$ in $D^{\perp_E}$ must lie in $\Hull_E(D)$.
\end{remark}

\subsection{Bookkeeping identities, and an erratum}

\begin{theorem}[bookkeeping]\label{thm:bookkeeping}
For every code produced by any of the three routes from $C\le\Rn$ with
$D=\varphi_n(C)$:
\begin{enumerate}[label=\textnormal{(\alph*)},leftmargin=*]
\item $k=\dim_{\F_q}D$ is even, and $\Nlen=4tn$ is divisible by $4$;
\item $K-c=\Nlen-2k\equiv0\pmod4$; in particular $K\equiv c\pmod4$;
\item $c+K=\Nlen-2h_E$;
\item $c\ge 2k-\Nlen$.
\end{enumerate}
\end{theorem}

\begin{proof}
(a) Under the Morita correspondence \eqref{eq:morita}, $D$ is the image of
$I_W$ with $W\le\Bq^2$, and $\dim_{\F_q}I_W=2\dim_{\F_q}W$; hence $k$ is even.
(b) $K-c=(\Nlen-k-h_E)-(k-h_E)=\Nlen-2k$, which is divisible by $4$ by (a)
and $\Nlen=4tn$. (c) is immediate from the two defining identities. (d) From
$c=k-h_E$ and $h_E\le\min(k,\Nlen-k)$ (as $h_E$ is the dimension of a subspace
of both $D$ and $D^{\perp_E}$) we get $c\ge k-(\Nlen-k)=2k-\Nlen$.
\end{proof}

\begin{remark}[erratum]
An earlier version of this project stated the identity as $c+K=\Nlen-k$. That
is wrong: it would force $h_E=k/2$, which is false in general, and it is
contradicted by $23$ of the $25$ rows of the earlier paper's own parameter
table. The correct identity is Theorem~\ref{thm:bookkeeping}(c). We record this
as an erratum rather than repair it silently; no \emph{parameter} computed in
the earlier work is affected, because $K$ and $c$ were always computed from
$k$ and $h_E$ directly.
\end{remark}

\subsection{Certification methodology}

Every distance reported in this paper is either \textbf{exact} or a
\textbf{proved lower bound}, and the two are never conflated. Three
independent methods are used, and every headline code is certified by all three.

\begin{description}
\item[M1 --- complete enumeration.] Enumerate the relevant classical difference
  set (the coset code of Theorem~\ref{thm:cosetdistance}, or the classical
  difference set of R2) completely and take the minimum weight. This is exact
  and complete whenever $q^{\dim}\le3\cdot10^6$.
\item[M2 --- increasing-weight support search.] A different algorithm: for each
  support of size $w=1,2,3,\dots$ in increasing order, solve the left-kernel
  problem asking whether the difference set contains a vector supported inside
  it. The first $w$ at which a solution exists is the distance. Because the
  algorithm differs structurally from M1, agreement of the two is a genuine
  cross-check.
\item[M3 --- ambient quantum search.] Search the ambient space
  $\F_q^{2\Nlen}$ in increasing symplectic weight for a vector $v$ with
  $v\in V^{\omegaperp}$ and $v\notin V$. This certifies the distance directly in
  the quantum language, without passing through
  Theorem~\ref{thm:cosetdistance}. It is complete only up to the weight reached,
  and is reported as a lower bound otherwise.
\end{description}

In addition, for every certified code: the stabilizer group is verified abelian
on an $\F_p$-basis (Theorem~\ref{thm:fpbasis}); $K$ and $c$ are recomputed from
$k$ and $h_E$; and --- crucially --- \emph{the certified code is the exact code
found by the search}, the image basis being carried through the whole pipeline.
An earlier version of this project certified a re-selected code with the same
invariants instead of the reported one; that defect was found and fixed, and the
current pipeline cannot reproduce it.

\subsection{Parameter analysis}

The computation covers $q\in\{4,8,9\}$ with $tn\le8$ (42 families), of which
$23$ are enumerated exhaustively, $4$ are counted but not enumerated, and $15$
are blocked by the component enumeration. Ten families that earlier work could
only sample are now exhaustive. In total $120\,721$ codes have computed
invariants and $5\,055$ have a computed distance, of which $4\,418$ are exact
and $637$ are proved lower bounds. Every one of the $5\,055$ satisfies the
entanglement-assisted Singleton and Hamming bounds, and the coset Singleton
bound \eqref{eq:cosetsingleton} (0 violations in each case).

\begin{proposition}[frontier, computational]\label{prop:frontier}
Among the $5\,055$ codes with a distance there are $368$ distinct tuples
$(\Nlen,K,d,c)$; $88$ of them are non-dominated and $280$ are dominated, where
$(N,K,d,c)$ dominates $(N',K',d',c')$ iff
$N\le N'$, $K\ge K'$, $d\ge d'$, $c\le c'$ with at least one strict
inequality, the comparison being made inside one alphabet.
\end{proposition}

\begin{remark}[what the frontier is and is not]
Proposition~\ref{prop:frontier} is a statement about the codes \emph{scanned},
not about all codes obtainable from the construction: the families that were
counted but not enumerated, and the fifteen that were blocked, are excluded, and
they are listed in the supplementary material. No tuple outside the scanned
families is claimed to be dominated.
\end{remark}

Table~\ref{tab:headline} lists the headline codes, all certified by the three
methods. Table~\ref{tab:cert} in Appendix~A lists the full certification table.

\begin{table}[t]
\centering
\caption{Headline codes. $S$-def is the entanglement-assisted Singleton defect
$\Nlen-K+c-2(d-1)$. All distances exact except the last row of the upper
block, which is a proved lower bound.}
\label{tab:headline}
\small
\begin{tabular}{@{}lllrrrrrr@{}}
\toprule
code & family $(q,t,n)$ & route & $k$ & $h_E$ & $K$ & $d$ & $c$ & $S$-def\\
\midrule
$[[24,8,4]]_9$     & $(9,1,6)$ & R2-SO & 8  & 8 & 8  & 4  & 0  & 10\\
$[[20,4,4]]_4$     & $(4,1,5)$ & R2-SO & 8  & 8 & 4  & 4  & 0  & 10\\
$[[20,4,4]]_4$     & $(4,1,5)$ & R2-DC & 12 & 8 & 4  & 4  & 0  & 10\\
$[[16,4,3]]_4$     & $(4,1,4)$ & R2-SO & 6  & 6 & 4  & 3  & 0  & 8\\
$[[16,4,3]]_4$     & $(4,1,4)$ & R2-DC & 10 & 6  & 4 & 3  & 0  & 8\\
$[[20,2,10;18]]_4$ & $(4,1,5)$ & R3    & 18 & 0 & 2  & 10 & 18 & 18\\
$[[24,2,10;18]]_9$ & $(9,1,6)$ & R3    & 20 & 2 & 2  & 10 & 18 & 22\\
$[[20,4,9;16]]_9$  & $(9,1,5)$ & R3    & 16 & 0 & 4  & 9  & 16 & 16\\
$[[24,4,9;20]]_4$  & $(4,1,6)$ & R3    & 20 & 0 & 4  & 9  & 20 & 24\\
$[[16,2,8;14]]_9$  & $(9,1,4)$ & R3    & 14 & 0 & 2  & 8  & 14 & 14\\
$[[16,2,7;10]]_9$  & $(9,1,4)$ & R3    & 12 & 2 & 2  & 7  & 10 & 12\\
$[[12,2,6;10]]_4$  & $(4,1,3)$ & R3    & 10 & 0 & 2  & 6  & 10 & 10\\
$[[24,4,6;20]]_4$  & $(4,2,3)$ & R3    & 20 & 0 & 4  & 6  & 20 & 30\\
$[[16,4,5;8]]_4$   & $(4,1,4)$ & R3    & 10 & 2 & 4  & 5  & 8  & 12\\
$[[24,10,4;2]]_9$  & $(9,1,6)$ & R3    & 8  & 6 & 10 & $\ge4$ & 2 & 10\\
\midrule
$[[24,6,2;2]]_9$   & $(9,2,3)$ & R3 & 10 & 8  & 6 & 2 & 2 & 18\\
$[[24,14,2]]_9$    & $(9,2,3)$ & R1 & 10 & --- & 14 & 2 & 0 & 8\\
$[[32,8,2;4]]_4$   & $(4,2,4)$ & R3 & 14 & 10 & 8 & 2 & 4 & 26\\
\bottomrule
\end{tabular}
\end{table}

The last three rows are the non-principal-hull witnesses of
Theorem~\ref{thm:nonprincipal}: their hulls have $\MU_{\Bq}=3$. They are
included for their structural significance, not for their parameters
($d=2$). The middle one of the three is a \emph{stabilizer} code produced by
route~R1, which is the only route available there because
Theorem~\ref{thm:obstruction} blocks the Euclidean one at $(q,t)=(9,2)$.

\subsection{Bounds satisfied}

\begin{proposition}[bounds]\label{prop:bounds}
All $5\,055$ codes with a computed distance satisfy
\begin{enumerate}[label=\textnormal{(\alph*)},leftmargin=*]
\item the entanglement-assisted Singleton bound $\Nlen-K+c\ge2(d-1)$;
\item the entanglement-assisted Hamming bound with effective dimension $K-c$:
  $\bigl(\sum_{i=0}^{t}\binom{\Nlen}{i}(q^2-1)^i\bigr)q^{K-c}\le q^{\Nlen}$
  with $t=\lfloor(d-1)/2\rfloor$;
\item the coset Singleton bound $K+d\le\Nlen+1$.
\end{enumerate}
A lower-bound distance is tested with its proved value.
\end{proposition}

\subsection{The entanglement budget}

Table~\ref{tab:budget} is the measured answer to the question behind the
criterion of Theorem~\ref{thm:feasibility}: what is the least entanglement at
which a given distance becomes available.

\begin{table}[t]
\centering
\caption{Minimum entanglement $c$ for each distance, over the whole search,
with the attaining code. The distance range is $d=2,\dots,8$: these are the
distances for which Part~4(c) of the Phase-7 log records a minimum-entanglement
attaining code. The two larger distances of the search, $d=9$ and $d=10$, are
attained only at $c=16$ and $c=18$ (Table~\ref{tab:headline}).}
\label{tab:budget}
\begin{tabular}{@{}clll@{}}
\toprule
$d$ & minimal $c$ & attaining code & status of $d$\\
\midrule
2 & 0  & $[[24,20,2;0]]_4$  & exact\\
3 & 0  & $[[24,12,3;0]]_9$  & exact\\
4 & 0  & $[[24,8,4;0]]_9$   & exact\\
5 & 2  & $[[24,6,5;2]]_9$   & proved lower bound\\
6 & 10 & $[[12,2,6;10]]_4$  & exact\\
7 & 10 & $[[16,2,7;10]]_9$  & exact\\
8 & 14 & $[[16,2,8;14]]_9$  & exact\\
\bottomrule
\end{tabular}
\end{table}

Reading Table~\ref{tab:budget} with Theorem~\ref{thm:cosetdistance}: $d=4$ is
available with \emph{no} entanglement at all ($[[24,8,4]]_9$, a stabilizer
code), while $d\ge5$ costs entanglement, and the cost grows steeply. With the
entanglement bounded by $c\le4$ the best attainable distance is $d=4$, attained
with $K=10$ and $c=2$ by $[[24,10,4;2]]_9$ --- a proved lower bound, and the
frontier point for the criterion.

\subsection{Worked example VI: a full certification}

% The optional argument is braced: LaTeX's [..] scanner stops at the first
% unbraced ']', and the quantum-code tuple contains ']]'.
\begin{example}[{certifying $[[20,2,10;18]]_4$}]\label{ex:certify}
Take the family $(q,t,n)=(4,1,5)$, so $\Nlen=4\cdot1\cdot5=20$. The search
finds a cyclic code $C\le\Rn$ whose Gray image $D=\varphi_5(C)\subseteq
\F_4^{20}$ has
\[
  k=\dim_{\F_4}D=18,\qquad h_E=\dim\bigl(D\cap D^{\perp_E}\bigr)=0 .
\]
Route R3 (Theorem~\ref{thm:R3}) then gives
\[
  \begin{gathered}
    c=k-h_E=18,\qquad K=\Nlen-k-h_E=20-18-0=2,\\
    [[20,2,d;18]]_4 .
  \end{gathered}
\]
The distance. By Theorem~\ref{thm:cosetdistance}, $d$ is the minimum weight of
the quotient code $\mathcal Q=D^{\perp_E}/\Hull_E(D)$; here $h_E=0$, so
$\Hull_E(D)=\{0\}$ and $\mathcal Q=D^{\perp_E}$, a linear $[20,2]_4$ code. Its
$4^2=16$ codewords are enumerated completely (method M1), and the minimum
weight found is $10$; the independent support search (M2) and the ambient
quantum search (M3, which certifies $d\ge3$ as a lower bound) agree with the
reported value. Hence $d=10$ exactly and
\[
  [[20,2,10;18]]_4 .
\]
The three bound checks:
\begin{enumerate}[label=\textnormal{(\roman*)},leftmargin=*]
\item entanglement-assisted Singleton:
  $\Nlen-K+c=20-2+18=36\ge2(d-1)=18$, with defect $18$;
\item entanglement-assisted Hamming with effective dimension $K-c=-16$:
  $\bigl(\sum_{i=0}^{4}\binom{20}{i}(4^2-1)^i\bigr)4^{-16}\le4^{20}$;
\item coset Singleton (Theorem~\ref{thm:cosetsingleton}):
  $K+d=2+10=12\le\Nlen+1=21$.
\end{enumerate}
For comparison, the entanglement-assisted Singleton bound alone only gives
$K\le\Nlen+c-2d+2=20$, so the coset Singleton bound is the binding one here
(Proposition~\ref{prop:newsbound}).

Finally, the bookkeeping identities of Theorem~\ref{thm:bookkeeping}:
$k=18$ is even; $K-c=2-18=-16=\Nlen-2k=20-36$ is divisible by $4$; and
$c+K=20=\Nlen-2h_E$, in accordance with the corrected identity --- the
earlier formula $c+K=\Nlen-k$ would give $20=2$, which is false.

No code found anywhere in the search has $d>10$; together with the
$[[24,2,10;18]]_9$ of Example~\ref{ex:routes}, this $[[20,2,10;18]]_4$ is
one of the two largest-distance codes of the search, and it is also the
most entanglement-hungry ($c=18$ out of $\Nlen=20$). That trade-off is
the content of Theorem~\ref{thm:feasibility}: $d\ge4$ together with
$K\ge2$ forces $2k\le\Nlen+c_0-K_0$, i.e. large $d$ and moderate $K$ cannot
both be had at small $c$.
\end{example}

% =====================================================================
\section{Comparison with existing quantum codes}
\label{sec:cmp}
% =====================================================================

\subsection{The terms of the comparison}

The comparison is with the tables of \cite{GrasslCodetables}, retrieved
2026-11-04 (page last changed 10.06.2024). Two limitations must be stated
before any number is quoted, and we state them where they bite.

\begin{enumerate}[label=\textnormal{(\arabic*)},leftmargin=*]
\item \textbf{The tables cover stabilizer codes only.} There is no
  entanglement-assisted table: the corresponding endpoint returns HTTP 404.
  Every entanglement-assisted code of this paper is therefore
  \emph{unbenchmarked}. This is a substantial blind spot, since the
  entanglement-assisted route is where the largest distances of this paper
  occur.
\item \textbf{Only $q=4$ and $q=9$ are tabulated.} For $q=8$ the site returns
  ``currently not available for q=8''. Hence the $(8,\cdot,\cdot)$ codes of this
  paper are unbenchmarked as well. (A third point of convention: the site prints
  the \emph{prime} $p$ in its code labels, so that $q=4$ appears as
  $[[N,K]]_2$ and $q=9$ as $[[N,K]]_3$, while the construction block says
  ``over GF($p^2$)''. An earlier phase of this project misread this and
  concluded, incorrectly, that $q=9$ was not covered; we record that as an
  erratum in the supplementary material and the corrected comparison is the one
  below.)
\end{enumerate}

\subsection{The comparison table}

\begin{table}[t]
\centering
\caption{Every stabilizer code ($c=0$) produced by the framework, compared
with \cite{GrasslCodetables}. ``Best known'' is the lower .. upper bound
pair for $[[N,K]]_q$.}
\label{tab:comparison}
\begin{tabular}{@{}lllc@{}}
\toprule
code & family $(q,t,n)$ & route & best known $[[N,K]]_q$ / gap\\
\midrule
$[[16,12,2]]_4$ & $(4,1,4)$ & R2-SO & $2..2$ --- meets\\
$[[16,12,2]]_4$ & $(4,1,4)$ & R2-DC & $2..2$ --- meets\\
$[[20,16,2]]_4$ & $(4,1,5)$ & R2-SO & $2..2$ --- meets\\
$[[20,16,2]]_4$ & $(4,1,5)$ & R2-DC & $2..2$ --- meets\\
$[[24,20,2]]_4$ & $(4,1,6)$ & R2-DC & $2..2$ --- meets\\[2pt]
$[[20,12,2]]_4$ & $(4,1,5)$ & R2-SO & $3..3$ --- 1 below\\
$[[20,12,2]]_4$ & $(4,1,5)$ & R2-DC & $3..3$ --- 1 below\\
$[[24,16,2]]_4$ & $(4,1,6)$ & R2-DC & $3..3$ --- 1 below\\[2pt]
$[[16,4,3]]_4$  & $(4,1,4)$ & R2-SO & $5..5$ --- 2 below\\
$[[16,4,3]]_4$  & $(4,1,4)$ & R2-DC & $5..5$ --- 2 below\\
$[[20,4,4]]_4$  & $(4,1,5)$ & R2-SO & $6..6$ --- 2 below\\
$[[20,4,4]]_4$  & $(4,1,5)$ & R2-DC & $6..6$ --- 2 below\\[2pt]
$[[16,8,3]]_9$  & $(9,1,4)$ & R2-DC & $4..4$ --- 1 below\\
$[[24,8,4]]_9$  & $(9,1,6)$ & R2-SO & $6..8$ --- 2 below (best known 6)\\
$[[24,12,3]]_9$ & $(9,1,6)$ & R2-DC & $5..6$ --- 2 below (best known 5)\\
\bottomrule
\end{tabular}
\end{table}

\subsection{Reading of the comparison}

Every stabilizer code constructed here is at best \emph{equal} to the best known
distance, and then only with $d=2$, where equality is uninformative. Every code
with $d\ge3$ is one or two below the best known value. \textbf{No record
parameter is claimed}, and none is implied: the codes are presented for their
origin, not for their position in the tables.

That said, the comparison does say something. The framework produces, at
$N=24$ over $\F_9$, a stabilizer code of distance $4$ where the best known is
$6$ (and where the upper bound is $8$): the construction is within a factor of
$1.5$ of the best known code at that length and dimension, and it does so from a
chain of entirely algebraic operations over a non-commutative ring. More
importantly, the \emph{entanglement-assisted} codes --- where this framework's
distances are largest ($d=10$ at $N=20$) --- cannot be compared at all, and the
absence of an entanglement-assisted table is a limitation of the literature
rather than of the construction.

\subsection{Novelty}

We do not claim that matrix rings, Galois hulls, Gray maps, quasi-cyclic images,
stabilizer codes or entanglement-assisted codes are new; each is attributed in
Section~\ref{sec:lit}. The novelty claim is the following, specific and proved.

\begin{enumerate}[label=\textnormal{(\arabic*)},leftmargin=*]
\item \textbf{The duality.} Over the non-commutative ring $R=M_2(A_t)$, the
  Galois twist must be $\theta_s=\sigma_s\circ(\cdot)^{\mathsf T}$, an
  \emph{anti}-automorphism. Over a commutative ring this degenerates to
  $\sigma_s$ and there is nothing to prove; here it is what makes the dual of a
  \emph{right} module a right module with the correct size
  (Theorems~\ref{thm:duality} and~\ref{lem:failure}).
\item \textbf{Hulls without principality, and the refutation.} The entire hull
  theory is developed for arbitrary right ideals, and the natural conjecture
  that Galois hulls are principal is \textbf{false}
  (Theorem~\ref{thm:nonprincipal}), with the phenomenon confined by the
  structural condition $t\ge2$ and $p\mid n$. This has no commutative analogue.
\item \textbf{The exact induced form.} The coefficient expansion transports the
  Galois pairing to $u^{[s]}(I_{4n}\otimes H_t)v^{\mathsf T}$
  (Theorem~\ref{thm:gram}) --- a $\widehat\sigma_s$-Hermitian form with a
  Hankel Gram matrix, not the Euclidean or plain Hermitian form unless $t=1$.
  Duality and hulls are preserved including on the non-principal hulls
  (Theorems~\ref{thm:grayduality} and~\ref{thm:grayhull}).
\item \textbf{Self-orthogonality without reflexivity.}
  Theorem~\ref{thm:hullso} is proved by a one-sided argument, so the whole
  construction is available in the general Galois regime.
\item \textbf{Two new bounds.} The coset-distance characterisation
  (Theorem~\ref{thm:cosetdistance}) and the coset Singleton bound
  (Theorem~\ref{thm:cosetsingleton}), the latter being neither the
  entanglement-assisted Singleton bound nor implied by it
  (Proposition~\ref{prop:newsbound}).
\item \textbf{Quantum codes from non-principal hulls.} The codes
  $[[24,6,2;2]]_9$ and $[[32,8,2;4]]_4$ of Table~\ref{tab:headline} come from
  hulls with $\MU_{\Bq}=3$. Their parameters are weak; their existence is the
  point, since it shows that no principality hypothesis is hidden anywhere in
  the chain.
\end{enumerate}

Relative to the closest prior work \cite{ZhangKongZheng2026QuantumGaloisHulls}
--- Galois hulls, a Gray map, and quantum codes over a finite non-chain ring ---
the differences are: their ring is \emph{commutative}, so no anti-automorphism
arises and no principality question can be asked; their Gray image has length
$tn$, whereas ours has length $4tn$ (Theorem~\ref{thm:nphys}), because
$\dim_{\F_q}M_2(A_t)=4t$; they determine hull \emph{dimensions}, whereas we
construct codes from hulls that are not principal and classify when
$\MU_{\Bq}\ge3$ can occur; and we prove the two bounds of item (5).

\subsection{On the earlier tuple $[[24,4,5;12]]_4$}

An earlier phase of this project reported $[[24,4,5;12]]_4$, obtained at
$(q,t,n)=(4,2,3)$ from a $300$-code sample. Exhaustive enumeration of that
family ($35\,937$ codes) shows that its best entanglement-assisted distance is
$6$, attained by $[[24,4,6;20]]_4$ ($k=20$, $h_E=0$, $c=20$, $K=4$): the
earlier value was a sampled lower bound, not an optimum. The tuple
$[[24,4,5;12]]_4$ is not reproduced by the exhaustive search and is dominated
(by $[[20,4,6;12]]_4$). We report this as a change of value obtained by a
stronger method, not as a retraction: nothing in the earlier derivation was
wrong, it was simply incomplete.

% =====================================================================
\section{Conclusion}
\label{sec:conclusion}
% =====================================================================

We have given a complete and self-contained treatment of the chain
\[
  \text{non-commutative matrix rings}
  \;\longrightarrow\;
  \text{cyclic codes}
  \;\longrightarrow\;
  \text{Galois duality and hulls}
  \;\longrightarrow\;
  \text{Gray maps}
  \;\longrightarrow\;
  \text{quantum codes}.
\]

The algebraic core is that over $R=M_2(A_t)$ the Galois twist must be the
anti-automorphism $\theta_s=\sigma_s\circ(\cdot)^{\mathsf T}$, so that the dual
of a right ideal is a right ideal with the correct dimension
(Theorems~\ref{thm:duality}, \ref{thm:dualformula}, \ref{lem:failure}). The hull
theory is developed for arbitrary --- in particular non-principal --- right
ideals, and the natural conjecture that Galois hulls are principal is false
(Theorem~\ref{thm:nonprincipal}), with the phenomenon confined to $t\ge2$ and
$p\mid n$.

The coefficient expansion transports the Galois pairing to the form
$u^{[s]}(I_{4n}\otimes H_t)v^{\mathsf T}$ (Theorem~\ref{thm:gram}); duality and
hulls are transported exactly, non-principal hulls included; and the image of a
Galois hull is self-orthogonal \emph{without} any reflexivity hypothesis
(Theorem~\ref{thm:hullso}), which is what makes the whole chain work in the
general Galois regime. A structural obstruction rules out the Euclidean
coordinate change in characteristic two with $t$ even
(Theorem~\ref{thm:obstruction}), and its hypotheses are exact.

On the quantum side, three routes (R1--R3) are available, their parameters and
distances are given with the hypotheses visible, and the distance is
characterised as the minimum weight of a quotient code
(Theorem~\ref{thm:cosetdistance}), yielding the coset Singleton bound
$K+d\le\Nlen+1$ (Theorem~\ref{thm:cosetsingleton}). A feasibility inequality
(Theorem~\ref{thm:feasibility}) describes the entanglement budget, and the
measured frontier is reported (Proposition~\ref{prop:frontier},
Table~\ref{tab:budget}).

The best codes obtained are $[[24,8,4]]_9$ among stabilizer codes and
$[[20,2,10;18]]_4$ and $[[24,2,10;18]]_9$ among entanglement-assisted codes.
\textbf{No record parameter is claimed.} The stabilizer codes are one or two
below the best known distances of \cite{GrasslCodetables}, and the
entanglement-assisted codes cannot be compared at all, since no
entanglement-assisted table exists. What is claimed is a new \emph{origin} for
quantum codes: non-commutative Galois hulls, transported exactly, with
non-principal hulls admitted and with two new bounds on the resulting
parameters.

\paragraph{Open problems.}
(i) Determine the Galois hull \emph{dimension} distribution in the wild
families, where the cyclic-code enumeration has resisted complete treatment for
$15$ families.
(ii) Find a closed or recursive description of the set of weights of the
quotient code of Theorem~\ref{thm:cosetdistance} in terms of the CRT data of
$C$.
(iii) Prove or refute the two conjectures stated at the end of
Section~\ref{sec:hulls}: $\Hull_s(C)\subseteq\sigma_{2s}(C)$, and
$\dim\Hull_s(C)=\dim\Hull_s(C^{\perp_s})$.
(vi) Decide the sharper form {\rm (O-3.1b)} of the refuted conjecture
(Remark~\ref{rem:o31b}): every witness currently known is a code that is
itself non-principal, so the question whether a \emph{principal} cyclic code
can have a non-principal Galois hull in the wild regime is open.
(iv) Determine whether the obstruction of Theorem~\ref{thm:obstruction} is the
only one: we have not proved that a congruence exists for all $t$ odd in
characteristic two, only that it exists for $t=1,3,5$ at $q=4,8$.
(v) Extend the comparison to $q=8$ and to entanglement-assisted codes, neither
of which is tabulated.


\appendix
% =====================================================================
\appendix
\section{Certification table and comparison data}
\label{app:tables}
% =====================================================================

\subsection{The certification table}

Table~\ref{tab:cert} is the complete output of the certification stage: all
$28$ certified codes, reproduced verbatim from the computation log
(\texttt{04\_COMPUTATION/phase7\_extended\_search\_log.txt}, Part~6). The
columns M1, M2, M3 record the three independent distance methods of
Section~\ref{sec:cert}: a dash means the method could not complete, and
``$>2$'' means the ambient quantum search (M3) reached weight $2$ without
finding a vector, i.e. it certifies $d\ge3$ as a \emph{lower} bound only.
The column ``exact'' says yes when all three methods that completed agree
with the reported value, the stabilizer is $\F_p$-abelian, and $K$ and $c$
recompute from $k$ and $h_E$.

\begin{table}[h]
\centering
\small
\caption{All $28$ certified codes (Part~6 of the Phase-7 log, verbatim).
$S$-def is the entanglement-assisted Singleton defect $N-K+c-2(d-1)$.}
\label{tab:cert}
\begin{tabular}{@{}llcrrrrrl@{}}
\toprule
code & family & route & $k$ & $K$ & $d$ & $c$
 & $S$-def & M1 / M2 / M3\\
\midrule
$[[24,8,4]]_9$   & $(9,1,6)$ & SO & 8  & 8  & 4  & 0  & 10 & -- / 4 / $>2$\\
$[[20,4,4]]_4$   & $(4,1,5)$ & SO & 8  & 4  & 4  & 0  & 10 & -- / 4 / $>2$\\
$[[20,4,4]]_4$   & $(4,1,5)$ & DC & 12 & 4  & 4  & 0  & 10 & -- / 4 / $>2$\\
$[[20,2,10;18]]_4$ & $(4,1,5)$ & EA & 18 & 2  & 10 & 18 & 18 & 10 / n/a / $>2$\\
$[[24,2,10;18]]_9$ & $(9,1,6)$ & EA & 20 & 2  & 10 & 18 & 22 & 10 / n/a / $>2$\\
$[[4,2,2;2]]_4$  & $(4,1,1)$ & EA & 2  & 2  & 2  & 2  & 2  & 2 / n/a / 2\\
$[[8,4,3;4]]_4$  & $(4,1,2)$ & EA & 4  & 4  & 3  & 4  & 4  & 3 / n/a / 3\\
$[[12,2,6;10]]_4$ & $(4,1,3)$ & EA & 10 & 2  & 6  & 10 & 10 & 6 / n/a / $>2$\\
$[[16,4,5;8]]_4$ & $(4,1,4)$ & EA & 10 & 4  & 5  & 8  & 12 & 5 / n/a / $>2$\\
$[[24,4,9;20]]_4$ & $(4,1,6)$ & EA & 20 & 4  & 9  & 20 & 24 & 9 / n/a / $>2$\\
$[[8,4,2;4]]_4$  & $(4,2,1)$ & EA & 4  & 4  & 2  & 4  & 6  & 2 / n/a / 2\\
$[[16,4,4;12]]_4$ & $(4,2,2)$ & EA & 12 & 4  & 4  & 12 & 18 & 4 / n/a / $>2$\\
$[[24,4,6;20]]_4$ & $(4,2,3)$ & EA & 20 & 4  & 6  & 20 & 30 & 6 / n/a / $>2$\\
$[[12,6,2;6]]_4$ & $(4,3,1)$ & EA & 6  & 6  & 2  & 6  & 10 & 2 / n/a / 2\\
$[[16,6,2;10]]_4$ & $(4,4,1)$ & EA & 10 & 6  & 2  & 10 & 18 & 2 / n/a / 2\\
$[[4,2,2;2]]_8$  & $(8,1,1)$ & EA & 2  & 2  & 2  & 2  & 2  & 2 / n/a / 2\\
$[[8,4,3;4]]_8$  & $(8,1,2)$ & EA & 4  & 4  & 3  & 4  & 4  & 3 / n/a / 3\\
$[[12,2,6;10]]_8$ & $(8,1,3)$ & EA & 10 & 2  & 6  & 10 & 10 & 6 / n/a / $>2$\\
$[[12,4,2;8]]_8$ & $(8,3,1)$ & EA & 8  & 4  & 2  & 8  & 14 & 2 / n/a / 2\\
$[[4,2,2;2]]_9$  & $(9,1,1)$ & EA & 2  & 2  & 2  & 2  & 2  & 2 / n/a / 2\\
$[[8,2,4;6]]_9$  & $(9,1,2)$ & EA & 6  & 2  & 4  & 6  & 6  & 4 / n/a / $>3$\\
$[[12,2,5;6]]_9$ & $(9,1,3)$ & EA & 8  & 2  & 5  & 6  & 8  & 5 / n/a / $>2$\\
$[[16,2,8;14]]_9$ & $(9,1,4)$ & EA & 14 & 2  & 8  & 14 & 14 & 8 / n/a / $>2$\\
$[[20,4,9;16]]_9$ & $(9,1,5)$ & EA & 16 & 4  & 9  & 16 & 16 & 9 / n/a / $>2$\\
$[[8,4,2;4]]_9$  & $(9,2,1)$ & EA & 4  & 4  & 2  & 4  & 6  & 2 / n/a / 2\\
$[[16,4,4;12]]_9$ & $(9,2,2)$ & EA & 12 & 4  & 4  & 12 & 18 & 4 / n/a / $>2$\\
$[[12,4,2;8]]_9$ & $(9,3,1)$ & EA & 8  & 4  & 2  & 8  & 14 & 2 / n/a / 2\\
$[[24,6,4;18]]_9$ & $(9,3,2)$ & EA & 18 & 6  & 4  & 18 & 30 & -- / 4 / $>2$\\
\bottomrule
\end{tabular}
\end{table}

\subsection{The non-principal witnesses}

The witnesses to Theorem~\ref{thm:nonprincipal} and the quantum codes they
produce, from Part~5 of the same log:

\begin{center}
\begin{tabular}{@{}ll@{}}
\toprule
witness & codes produced\\\midrule
$(9,2,3)$, $s=0$ & R1: $[[24,14,2]]_9$ (abelian); R3: $[[24,6,2;2]]_9$\\
$(9,2,3)$, $s=1$ & R1: $[[24,14,2]]_9$ (abelian); R3: $[[24,6,2;2]]_9$\\
$(4,2,4)$, $s=0$ & R3: $[[32,8,2;4]]_4$\\
$(4,2,4)$, $s=1$ & R3: $[[32,8,2;4]]_4$\\
\bottomrule
\end{tabular}
\end{center}

In all four cases $\MU_{\Bq}(W)=3$ and $\MU_{\Bq}(\Hull_s(W))=3$. The two
$(9,2,3)$ witnesses additionally have $\Hull_s(W)=W$, i.e. the code equals its
own Galois hull.

\subsection{Best known distances used in Section~\ref{sec:cmp}}

Table~\ref{tab:grassl} is the raw data retrieved from
\cite{GrasslCodetables} on 2026-11-04, in the form lower..upper bound for
$[[N,K]]_q$. The entry ``not available'' is the site's own response for
$q=8$.

\begin{table}[h]
\centering
\caption{Best known stabilizer distances, codetables.de, retrieved
2026-11-04. The site prints the \emph{prime} $p$ in its code labels, so
$q=4$ appears as $[[N,K]]_2$ and $q=9$ as $[[N,K]]_3$.}
\label{tab:grassl}
\begin{tabular}{@{}cc@{\qquad}cc@{}}\toprule
$[[N,K]]_q$ & lower..upper & $[[N,K]]_q$ & lower..upper\\\midrule
$[[16,4]]_4$   & 5..5  & $[[8,4]]_9$   & 3..3\\
$[[16,12]]_4$  & 2..2  & $[[16,8]]_9$  & 4..4\\
$[[20,4]]_4$   & 6..6  & $[[24,8]]_9$  & 6..8\\
$[[20,12]]_4$  & 3..3  & $[[24,10]]_9$ & 6..7\\
$[[20,16]]_4$  & 2..2  & $[[24,12]]_9$ & 5..6\\
$[[24,16]]_4$  & 3..3  & $[[8,4]]_8$   & not available\\
$[[24,20]]_4$  & 2..2  & $[[12,2]]_8$  & not available\\
\bottomrule
\end{tabular}
\end{table}

% =====================================================================
\section{Reproducibility}
\label{app:repro}
% =====================================================================

Every number in this paper that is not proved in the text is produced by a
script in the accompanying supplementary package, and every such number is
attributable to a line of a committed log. This appendix records what the
package contains, how to run it, and which numbers come from where.

\subsection{Package layout}

\begin{verbatim}
PAPER/supplementary/
  README.md                  index of the package
  ENVIRONMENT.md             exact environment used
  MANIFEST.sha256            SHA-256 of every file in the package
  scripts/                   the computation scripts (copies of the
                             committed 04_COMPUTATION/ scripts)
  logs/                      the committed logs the numbers are read from
  generator_matrices/        generator matrices of the headline codes,
                             one .json per code
  parameter_tables/          machine-readable versions of the tables of
                             this paper (CSV)
  CERTIFICATION_PROCEDURE.md what "certified" means here, and how to
                             re-certify a code from scratch
\end{verbatim}

The scripts are byte-identical copies of the scripts committed in
\texttt{04\_COMPUTATION/} of the project repository; the logs are the
committed logs. Nothing was re-run to produce the package, so no script hash
is invalidated.

\subsection{Environment}

The computation is pure Python~3 standard library. No third-party package
(numpy, sympy, GAP, Magma) is used or required; arithmetic in $\F_q$ and in
$A_t$, $R$, $\Rn$ is implemented in \texttt{matrix\_ring\_lib.py} and is
elementary by design, so that every step can be inspected. See
\texttt{ENVIRONMENT.md} for the interpreter version and platform of the
recorded run.

\subsection{Randomness}

The search uses pseudo-random sampling for the families that are too large to
enumerate. The seed is fixed (\texttt{SEED = 20261210}) and every random
choice is drawn from a single \texttt{random.Random(SEED)} instance in a
fixed order of families, so a re-run reproduces every line of the log except
the per-family elapsed-time annotations and the peak-memory line, which are
wall-clock and resource readings and take no part in any computation.

The principal results of this paper do not depend on the sampling: the
theorems are proved, and the headline codes come from families that are
enumerated exhaustively.

\subsection{Where the numbers come from}

\begin{itemize}[leftmargin=*]
\item Theorems~\ref{thm:duality}--\ref{thm:parameters},
  \ref{thm:cosetdistance}, \ref{thm:cosetsingleton},
  \ref{thm:feasibility}, \ref{thm:bookkeeping}, \ref{thm:gram},
  \ref{thm:obstruction}, \ref{thm:hullso}, \ref{thm:fpbasis},
  \ref{thm:nonprincipal}: proved in the text; no computation enters.
\item Table~\ref{tab:witnesses} (the $\MU_{\Bq}=3$ witnesses) and the four
  witness codes of Appendix~\ref{app:tables}: Part~5 of
  \texttt{phase7\_extended\_search\_log.txt}, produced by
  \texttt{phase7\_extended\_search.py}.
\item Table~\ref{tab:cert} (all 28 certified codes): Part~6 of the same log.
\item Table~\ref{tab:headline}: Part~4(d) (Pareto frontier), Part~4(e) (best
  per route) and Part~6 of the same log.
\item Table~\ref{tab:budget} (minimum entanglement per distance): Part~4(c)
  of the same log.
\item Table~\ref{tab:comparison} and Table~\ref{tab:grassl}: Part~8 of the
  same log for the $q=4$ rows, and the pages of \cite{GrasslCodetables}
  retrieved on 2026-11-04 for the $q=9$ rows and for the bounds themselves.
\item The counts $120\,721$, $5\,055$, $4\,418$, $637$, $368$, $88$, $280$,
  $25$ families scanned, $42$ families in the census ($23$ exhaustive, $4$
  counted-only, $15$ blocked): the completion report and Part~2 of the same
  log.
\item The bookkeeping identities of Theorem~\ref{thm:bookkeeping} are
  proved, and check A1 of the log verifies them on all $5\,055$ codes
  ($0$ violations). The coset Singleton bound is check H1 ($0$ violations).
  The Singleton and Hamming bounds are check G1 ($0$ violations).
\end{itemize}

\subsection{How to reproduce}

\begin{verbatim}
  cd PAPER/supplementary/scripts
  python3 phase7_extended_search.py > out.txt 2>&1
  diff <(grep -v '\[[0-9]*s\]' out.txt) \
       <(grep -v '\[[0-9]*s\]' ../logs/phase7_extended_search_log.txt)
\end{verbatim}

The run takes about $9$ minutes and about $370$~MB of resident memory on the
recorded platform. The \texttt{grep -v} filter removes the elapsed-time
annotations; the peak-memory line is the only other difference.

The earlier phases are reproduced the same way from their own scripts; see
\texttt{04\_COMPUTATION/README.md} in the repository for the dependency
order (\texttt{matrix\_ring\_lib.py} first, then the phase scripts).

\subsection{What is \emph{not} reproducible from the package}

The $q=9$ bounds of Table~\ref{tab:grassl} were read from a web page that
changes; the retrieved values are dated 2026-11-04 and are quoted verbatim,
and the page reported ``Last change: 10.06.2024''. They are external data,
not an output of this computation.

\subsection{Verification of the package itself}

\begin{verbatim}
  cd PAPER/supplementary
  sha256sum -c MANIFEST.sha256
\end{verbatim}

Every file listed in the manifest is checked against its recorded hash; see
\texttt{PAPER/SHA256\_MANIFEST.md} for the same check over the manuscript and
the whole package.


% ------------------------------ bibliography --------------------------
\bibliographystyle{plain}
\bibliography{references}

\end{document}
